\documentclass[11pt,a4paper,twoside]{article}
\usepackage[utf8]{inputenc}
\usepackage[english]{babel}
\usepackage{amsmath,amsfonts,amsthm,amssymb,mathtools}

\usepackage{xcolor}
\usepackage{tabularx}
\usepackage{tikz-cd}
\usepackage{graphicx}
\usepackage{epigraph}
\usepackage[hidelinks]{hyperref}
\hypersetup{
    colorlinks,
    linkcolor={red!50!black},
    citecolor={blue!50!black},
    urlcolor={blue!80!black}
}

\newtheorem{theorem}{Theorem}[section]
\newtheorem{corollary}[theorem]{Corollary}
\newtheorem{lemma}[theorem]{Lemma}
\newtheorem{proposition}[theorem]{Proposition}

\theoremstyle{definition}
\newtheorem{definition}[theorem]{Definition}

\theoremstyle{remark}
\newtheorem{remark}[theorem]{Remark}

\newtheorem{openquestion}[theorem]{Open Question}

\newtheoremstyle{example}
  {\topsep}   % ABOVESPACE
  {\topsep}   % BELOWSPACE
  {\normalfont}  % BODYFONT
  {0pt}       % INDENT (empty value is the same as 0pt)
  {\bfseries} % HEADFONT
  {.}         % HEADPUNCT
  {\newline} % HEADSPACE
  {}          % CUSTOM-HEAD-SPEC
\theoremstyle{example}
\newtheorem{example}[theorem]{Example}

\newcommand{\pp}[1]{\mathbb{P}[#1]}
\newcommand{\indf}[1]{\mathbf{1}_{\{#1\}}}
\newcommand{\qq}[1]{\mathbb{Q}[#1]}
\newcommand{\stt}{\text{ s.t }}
\newcommand{\extreals}{\overline{\mathbb{R}}}
\newcommand{\extrm}{\overline{\mathcal{R}}}

\newcommand{\drmrm}{\overline{\mathcal{R}_{\text{DRM}}}}
\newcommand{\extrmcv}{\overline{\mathcal{R}_{\text{cv}}}}
\newcommand{\extrmch}{\overline{\mathcal{R}_{\text{ch}}}}

\newcommand{\extrmloc}{\overline{\mathcal{R}_{\text{loc}}}}


\newcommand{\convh}[1]{\text{ch}\left ( #1 \right)}
\newcommand{\cconvh}[1]{\overline{\text{ch}}\left ( #1 \right)}
\newcommand{\cenv}{\text{cv}}
\newcommand{\ccenv}{\overline{\text{cv}}}

%\newcommand{\comcore}[1]{\text{ComCore} \left ( #1 \right )}
%\newcommand{\comconvh}[1]{\text{ComCH} \left ( #1 \right )}
\newcommand{\comint}[1]{\text{ComInt} \left ( #1 \right )}
\newcommand{\comhull}[1]{\text{ComHull} \left ( #1 \right )}
\newcommand{\maxhs}[2]{\text{MaxHS}_{#1} \left ( #2 \right )}
\newcommand{\minhs}[2]{\text{MinHS}_{#1} \left ( #2 \right )}

\newcommand{\dextrm}[2]{d_{\extrm} \left (#1, #2 \right)}
\newcommand{\acceptidealb}[1]{\overline{R}\left (#1 \right)}
\newcommand{\acceptideal}[1]{R \left (#1 \right)}
\newcommand{\accidbloc}[1]{\overline{R_{\text{loc}}}\left (#1 \right)}
\newcommand{\accidb}[2]{\overline{R}_{#1}\left (#2 \right)}
%\newcommand{\accidbl}{R \left (#1 \right)}

\newcommand{\rejectidb}[1]{\overline{F}\left (#1 \right)}
\newcommand{\rejidbloc}[1]{\overline{F_{\text{loc}}}\left (#1 \right)}
\newcommand{\rejectid}[1]{F\left (#1 \right)}
\newcommand{\rejidb}[2]{\overline{F}_{#1}\left (#2 \right)}
%\newcommand{\rejidbl}{\overline{F}_{#2}\left (#1 \right)}

\newcommand{\riskidealb}[1]{\overline{A}\left (#1 \right)}
\newcommand{\riskideal}[1]{A \left (#1 \right)}


\newcommand{\evalrmf}[1]{\textnormal{EVF}\left( #1 \right)}
\newcommand{\evalrmr}[1]{\textnormal{EVR}\left( #1 \right)}

\newcommand{\pmom}{\rho_{\text{mom}}}
\newcommand{\ppop}{\rho_{\text{pop}}}
\newcommand{\expv}[1]{\mathbb{E} \left [#1 \right]}
\newcommand{\pexpv}[2]{\mathbb{E}_{#1} \left [#2 \right]}
\newcommand{\prob}[1]{\mathbb{P} \left [#1 \right]}
\newcommand{\var}[1]{\text{Var}[#1]}
\newcommand{\floor}[1]{\lfloor #1 \rfloor}
\newcommand{\ind}{\perp\!\!\!\!\perp} 
\newcommand{\spec}[1]{\text{Spec } #1} 
\newcommand{\aff}[2]{\mathbb{A}^{#1}_{#2}}
\newcommand{\norm}[1]{\left|\left| #1 \right| \right|}
\newcommand{\cl}[1]{\text{cl}\left(#1\right)}
\newcommand{\innerpr}[2]{\langle {#1}, {#2} \rangle}
\def\resf#1{\kappa(\mathfrak{#1})}
\def\code#1{\texttt{#1}}

\newcommand{\idl}[1]{\text{Idl}\left (#1 \right)}
\newcommand{\upuparr}{ {\uparrow \uparrow }}
\newcommand{\dodoarr}{ {\downarrow \downarrow }}

\def\esssup{\textnormal{ess sup }}
\def\essinf{\textnormal{ess inf }}

\DeclareMathOperator*{\argmax}{arg\,max}
\DeclareMathOperator*{\argmin}{arg\,min}

\newcommand{\opencat}{\textbf{Open}}
\newcommand{\quotlat}[1]{\textbf{QLat}\left (#1 \right)}
\newcommand{\presh}{\mathcal{F}^{\text{pre}}}
\newcommand{\flt}[1]{\textnormal{Flt}({#1})}
\newcommand{\satflt}[2]{\textnormal{SatFlt}_{#1}({#2})}

\newcounter{coneofringscounter}


\title{Sheafs Of Risk Measure Lattices}
\author{N.J. Lohmann, 10064840}
\begin{document}
\maketitle
%\tableofcontents
\begin{abstract}
This paper presents a study of the local behavior of monetary risk measures using their natural lattice structure
and their geometric interpretation of operations on their acceptance sets. In particular, this allows to define
a sheaf with sections of quotient lattices of the risk measure lattice and to study their behavior locally.
While it is a fortunate development that this is conceptually possible it will also be highlighted, how this approach
is obstructed to study aspects that are of typical interest in the theory of risk measures. This is exemplified
with a short overview of the traditional local study of a sublattice, namely the one of convex risk measures.
Nevertheless, this is a novel take to apply the abstract theory of algebraic geometry in the context of financial mathematics
and also gives another category, the one of distributive lattices, to which sheaves can map a topological space.
To put the work in the context, it is an extension of the analysis conducted by the author in his master thesis, in
which general approaches via algebraic geometry to risk measure theory are treated.
\end{abstract}

\newpage
\begin{section}{Sheaves}
Sheaves are ubiquitous in mathematics and even though the reader might not had the opportunity
to apply them knowingly they can be sure to have used their structure implicitly in different categories.
From historic perspective, they originate in the domain of algebraic topology, where their ability to capture
consistenly the local behavior of objects such as continuous functions allowed an accessable
treatment of rather involved topological spaces. Over the course of time one of the greatest 
extensions of this \textit{tracking idea} has been achieved by Alexander Grothendieck, who used
them in his work EGA \cite{egatransl} in their specialized version as schemes to generalize algebraic geometry to the spectrum 
of rings, unifying geometric perspectives from the polynomial and number theoretic settings.
The sheaves in this work have a similar spirit, so without further ado:

\begin{definition}[Presheaves and Sheaves]
    Given a topological space \( X \), a sheaf is a structure s.t. 
    for all open \( U \subseteq X \), there is a map of sets \( \mathcal{O} : U \to \mathcal{O}(U) \) 
    and for all open \( V \subseteq U \) with their respective inclusion \( \iota : V \to U  \),
    corresponds a map \( r_{U,V} : \mathcal{O}(U) \to \mathcal{O}(V) \).
    The \( r \) stands for restriction and we use the notation \( f|_V := r_{U, V}(f) \) for \( f \in \mathcal{O}(U) \). Additionally 
    the maps should satisfy:
    \begin{enumerate}
        \item \( r_{U, U} = \text{id}_{\mathcal{O}(U)} \)
        \item If \( U \subseteq W \) then \( r_{U, V} \circ r_{W, U} = r_{W, V} \).
        \item (Locality): For a cover \( \bigcup_i U_i = U \) and  \( f, g \in \mathcal{O}(U) \), s.t. \( f|_{U_i} = g|_{U_i} \) 
            for all \( i \), then \( f = g \).
        \item (Gluing): For a cover \( \bigcup_i U_i = U \) and a collection \( \{ f_i \}_i,\ f_i \in \mathcal{O}(U_i) \), 
           s.t. \( f_i|_{U_i \cap U_j} = f_j|_{U_i \cap U_j} \), then there exists \( f \in \mathcal{O}(U) \)
           s.t. \( f|_{U_i} = f_i \).
    \end{enumerate}
    A structure which only satisfies the first two properties is a presheaf. If it also satisfies locality it is separated.
\end{definition} 
%The quite general definition hints already at the fact that sheaves are ubiquitous in mathematics, e.g. smooth functions over a manifold
%form a sheaf, the sheaf of Schwartz distributions over a space and there is even a sheaf of measures on Borel spaces \cite{sheafmeasures}.
%In the affine case here, there is a particular choice of sheaf, which got its own name. For this we also need to define morphisms.

\begin{example} \ref{algstructuresheaf}
    Consider the affine space \( \aff{n}{\mathbb{K}} = \mathbb{K}^n\) over some field \( \mathbb{K} \) 
    and the associated commutative ring of multivariate polynomials \( \mathbb{K}[\mathbf{X}] \).
    One can define the Zariski topology (a \( T_1 \)-topology) over \( \aff{n}{\mathbb{K}} \), whose subbasis of closed sets is induced through 
    \( J \subseteq \mathbb{K}[\mathbf{X}] \) via their vanishing sets, that is 
    \[ V(J) = \{\, a \in \mathbb{K}^n : f(a) = 0 \text{ for all } f \in J \,\}. \]
    Clearly, one also obtains a subbasis of open sets by taking the complements of vanishing sets (denoted by \( D(J) \)
    traditionally), which are points over which collection of polynomials do not vanish.
    Define then for an open set \( U \subset \aff{n}{\mathbb{K}} \) the set of \( \mathcal{O}(U) \) regular functions over \( U \), that is 
    \( f : U \to \mathbb{K} \), s.t. for every \( x \in U \) there is some open neighborhood \( V \subset U \) of it on which it holds
    \( f|_{V} \neq 0 \) and \( f(x) = \frac{h(x)}{g(x)} \) for \( h, g \in \mathbb{K}[\mathbf{X}] \).
    This mapping \( \mathcal{O} : \text{Open}(\aff{n}{\mathbb{K}}) \to \mathcal{O}(U) \) defines in fact a sheaf over \( \aff{n}{\mathbb{K}} \)
    with the restriction being the standard domain restriction of functions as the reader can verify. Moreover, the 
    \( \mathcal{O}(U) \) are also rings, since adding or multiplying two non-vanishing functions \( f_1, f_2 \) yields another one 
    over the intersection \( U_1 \cap U_2 \) of their domains, turning \( \aff{n}{\mathbb{K}} \) into a locally ringed space.
\end{example}

This example hopefully illuminates how the construction of the sheaf of smooth functions over a smooth manifold would look like.
As is apparent from the Definition, the local behavior of functions can be tracked through the restriction morphism. If we continue this 
zoom into smaller and smaller open sets indefinitely and end up at a point (which is more than often not an open set in the topology of the underlying space \( X \)) the result is often something similar to an evaluation. The formal construction of this idea requires the colimit 
from category theory, which can be found in the Appendix \ref{defcolimit}, the main challenge is to properly handle that it doesn't matter
from which open set we initialize our restrictions to the point. This idea has its own name for sheaves: 

\begin{definition}
    The stalk at a point \( x \in X \) is the colimit of the sheaf sections \( \mathcal{O}(U) \) of the open neighborhoods \( U \) of the point,
    denoted by
    \[
        \mathcal{O}_{X,x} = \varinjlim_{U\ni x} \mathcal{F}(U).
    \]
\end{definition}
\begin{example}
    For the sheaf of regular functions over \( \aff{n}{\mathbb{K}} \) the stalk at a point \( x \) is given by the ring 
    of rational polynomial functions \( f = \frac{g}{h} \) with \( h(x) \neq 0 \).
\end{example}
\end{section}

Everything is a matter of perspective and there is rarely one canonical choice over some space and its 
topology. Already for smooth manifolds we can assign for each \( k \in \mathbb{N} \) to an open set
the set of \( k \)-differentiable functions over it, giving a plethora of sheaves. It can be useful to inspect
how this different sheaves relate to each other as long as there is a consistent way of translating between their sections
which is made precise by the following idea.
\begin{definition}[Morphism of sheaves] \label{sheafmorphism}
    A morphism \(\varphi : \mathcal{F} \to \mathcal{G}\) of sheaves on \(X\) is a rule which assigns to each open \(U \subseteq X\) 
    a map of sets \(\varphi : \mathcal{F}(U) \to \mathcal{G}(U)\) compatible with restriction maps, i.e., whenever \(V \subseteq U \subseteq X\) are open 
    the following diagram commutes: 
    \[
        \begin{tikzcd}
            \mathcal{F}(U) \arrow[r, "\varphi"] \arrow[d, "r_{U,V}"] & \mathcal{G}(U) \arrow[d, "r_{U,V}"] \\
            \mathcal{F}(V) \arrow[r, "\varphi"] & \mathcal{G}(V)
        \end{tikzcd}
    \]
    An isomorphism of sheaves is a morphism \( \varphi : \mathcal{F} \to \mathcal{G} \), s.t there exists a morphism 
    \( \varphi^{-1} : \mathcal{G} \to \mathcal{F} \), satisfying 
    \begin{align*}
        \varphi^{-1} \circ \varphi &= \text{id}_{(X, \mathcal{F})} \\
    \varphi \circ \varphi^{-1} &= \text{id}_{(Y, \mathcal{G})} 
    \end{align*}
\end{definition}
Note that if there exists a homeomorphism \( \pi : X \to Y \) between topological spaces, we can
also define morphisms of sheaves with different underlying spaces, namely if \( (X, \mathcal{F}) \) and \( (Y, \mathcal{G}) \)
are sheaves, then for a open sets \( U \subseteq X, V \subseteq Y \) one could construct maps 
\( \varphi : \mathcal{F}(U) \to \mathcal{G}(\pi(U)) \), \( \varphi' : \mathcal{F}(\pi^{-1}(V)) \to \mathcal{G}(V)  \)
and if for all \( U \) (respectively \( V \)) these maps satisfy the definition of \ref{sheafmorphism},
we get a morphism between sheaves with different underlying space. Unfortunately, these can't be the same
(simply because \( X \) and \( Y \) are not the same space), but there's a one to one correspondence
(due the adjointness of the preimage sheaf and pushfoward sheaf functor). In case that \( \pi \)
is just a continuous map, one refers typically to morphisms of sheaves to the direction \( \varphi' : \mathcal{F}(\pi^{-1}(V)) \to \mathcal{G}(V) \) 
(so it is a morphism of sheaves on \( Y \)).

\newpage
\begin{section}{Risk Measures}
    \begin{subsection}{Properties, Examples and Obstructions}
    Risk measures are now a standard tool for evaluation in actuarial science as well as in financial mathematics as a whole. 
    One of the most basic risk measures,
    \[ 
        \rho(X) = \mathbb{E}[X] 
    \] 
    dates back as early as 1709 to Nicolaus I Bernoulli’s PhD. thesis \cite{bernoulli1709dissertatio}, in which he proposed it as the fair price for insurance contracts.
    % https://pmc.ncbi.nlm.nih.gov/articles/PMC11928225/
    % Roman Law on the Just Price in Nicolaus Bernoulli’s Mathematics
    Remarkably, only a few years later, Bernoulli also suggested an adjustment based on individual preferences—a concept we would now recognize as a precursor to utility theory.
    It took several centuries, however, until Markowitz’s portfolio theory \cite{markowitzportfolio}, which introduced variance as a measure of risk, reignited the theory and 
    led to the early development of downside risk measures such as lower partial moments. Specifically, in 1959, Markowitz (Chapter 9 \cite{markowitzsemivariance}) proposed
    \[
        \mathbb{E}[(|\mathbb{E}[X] - X|\mathbf{1}_{ \{\mathbb{E}[X] - X > 0 \} })^k] 
    \] 
    Over the following decades, various risk measures flourished, mostly under the term "premium principle", 
    % A General Class of Three-Parameter Risk Measures Bernell K. Stone
    of which special mentioning deserves the value-at-risk 
    \[ 
        \rho(X) = \text{VaR}_{\lambda}(X) = \inf \{ m \in \mathbb{R} | \pp{X + m < 0} \leq \lambda \}
    \]
    and the expected shortfall, or average value at risk ,
    \[
        \rho(X) = \text{AVR}_{\lambda}(X) = \frac{1}{\lambda} \int_{0}^\lambda \text{VaR}_{\alpha}(X) d\alpha
    \]
    until in 2011 Föllmer and Schied \cite{foellmerschied} put the theory on rigorous axiomatic foundations.
    % Risk Measures and Theories of Choice Tsanakas
    While an algebraic approach to portfolio theory is already well developed—since it is inherently an optimization problem on an intersection of hyperplanes and a conic—
    the algebraic aspect of risk measures still offers ample room for development, to the best of the author’s knowledge. 
    \newline

    We first give some basic definitions for risk measures.
    Let \( \mathcal{X} = \mathcal{L}^p(\Omega, \mathcal{F}, \mathbb{P}) \) or subspace thereof, 
    for some suitable choice of \( p \).
    \begin{definition} \label{defriskmeasure}
    A mapping \( \rho : \mathcal{X} \to \mathbb{R} \) is called a monetary risk measure if it satisfies
    the following conditions for all \( X, Y \in \mathcal{X} \).
        \begin{enumerate}
            \item Monotonicity/Antitone: If \( X \leq Y\), then \(\rho(X) \geq \rho(Y) \).
            \item Cash invariance: If \( m \in \mathbb{R} \), then \( \rho(X + m) = \rho(X) - m \).
        \end{enumerate}
        It is called normalized if \( \rho(0) = 0 \), convex if \( \rho \) is a convex functional, and coherent if it is in addition positively homogenous,
        that is \( \rho(\lambda X) = \lambda \rho(X),\ \lambda \in \mathbb{R}_{\geq 0} \).
    \end{definition}
    \begin{remark}
    If we say risk measure, we mean monetary risk measure, if not otherwise stated. The concept of coherence could be referred to as a sublinear 
    risk measure in functional analytic terminology.
    \end{remark}
    \begin{definition}
        A risk measure \( \rho \) is
        \begin{enumerate}
            \item sensitive, if \( \rho(-X) > \rho(0) \) for \( X \in L^\infty_{+}(\Omega, \mathcal{F}, \mathbb{P}) \) and \( \pp{X > 0} > 0 \),
            \item law-invariant, if \( \rho(X) = \rho(Y) \) for \( X \sim Y\) under \( \mathbb{P} \),
            \item continuous from below, if \( X_n \nearrow X \implies \rho(X_n) \searrow \rho(X) \),
            \item continuous from above, if \( X_n \searrow X \implies \rho(X_n) \nearrow \rho(X) \).
        \end{enumerate}
    \end{definition}

    \begin{definition}[Linear cones]
    A set \( S \) in a real vector space \( \mathcal{X} \) is a cone if \( x \in S \implies \lambda x \in S \) for \( \lambda \in \mathbb{R}_{>0} \).
    \end{definition}
    \begin{remark}
       To differentiate cones as in the previous from the ones such as in Definition \ref{coneofrings}
       we call the first linear cones and the latter algebraic cones.
    \end{remark}
    \begin{definition}
        The acceptance set of a risk measure \( \rho \) is 
        \[ \mathcal{A}_\rho = \{ X \in \mathcal{X} | \rho(X) \leq 0 \}. \]
    \end{definition}
    \begin{proposition} \label{acceptancesetpropisriskprop}
        Assume \( \mathcal{A}_{\rho} \neq \emptyset \), it holds:
        \begin{enumerate}
            \item \( \rho(X) := \inf \{ m \in \mathbb{R} | m + X \in \mathcal{A}_{\rho} \}   \) 
            \item \( \rho \) is a convex risk measure if and only if \( \mathcal{A}_{\rho} \) is convex.
            \item \( \rho \) is positive homogenous if and only if \( \mathcal{A}_{\rho} \) is a cone.
            \item \( \rho \) is coherent if and only if \( \mathcal{A}_{\rho} \) is a convex cone 
        \end{enumerate}
    \end{proposition}
    \begin{proof}
        \cite{foellmerschied} Proposition 4.6
    \end{proof}
    \end{subsection}

\end{section}
\newpage
\begin{section}{Algebraic Structures On Risk Measures}
    \begin{subsection}{Lattices}
    In this section we will provide the preliminaries from the theory of lattices for the sheaf construction
    on \( \mathcal{X} = L^\infty(\Omega, \mathcal{F}, \mathbb{P}) \), such as their definition, properties
    such as completeness and distributivity, which allow for a well-defined theory of ideals. Further a impression 
    of the geometric correspondence to acceptance sets is given, whose extension is explored in the subsequent sections.
    For the full proofs the reader may refer to the thesis \cite{mymaster}, which was the initial motivation of this work. 

    \begin{definition}
        A semilattice \( (S, \otimes ) \) is a set \( S \) with a binary operation which is 
        commutative, associative, and idempotent. It is called bounded if there's an element 
        \( e \in S \) s.t. \( x \otimes e = x \) for all \( x \in S \).
        Given two semilattices \( (L, \wedge), (L, \vee) \) on a set \( S \), whose operations
        satisfy the absorption laws:
        \begin{align*}
            a \vee (a \wedge b) = a  \\
            a \wedge (a \vee b) = a,
        \end{align*}
        the resulting structure \( (L, \wedge, \vee) \) is called a lattice and the operations \( \wedge \), \( \vee \) join and meet respectively.
        If both semilattices are bounded, the corresponding elements are denoted by \( \bot,\top \) respectively.
        A lattice is distributive if it in addition satisfies
        \begin{align*}
         a \vee (b \wedge c) = (a \wedge b) \vee (a \wedge c)  \\
         a \wedge (b \vee c) = (a \vee b) \wedge (a \vee c)  \\
        \end{align*}
    \end{definition}

    \begin{lemma}
        In a lattice \( (L, \wedge, \vee) \) for an element \( a \in L \) the sets
        \begin{align*}
            (a)^\uparrow &= \{ b \in L | a \leq b\}   \quad (a)^\downarrow = \{ b \in L | b \leq a\}  \\
            (a)^{\upuparr} &= \{ b \in L | a < b\} \quad (a)^{\dodoarr} = \{ b \in L | b < a\}  
        \end{align*}
        are semilattices, which we call the (strict) upward and (strict) downward of "\( a \)" respectively.
        If \( L \) is bounded, so are the semilattices.
    \end{lemma}

    \begin{proposition}
        There exists a unique partial order on a semilattice and for a lattice these orders agree.
    \end{proposition}
    \begin{proof}
        Define \( a \leq b \iff a \vee b = b \) and for lattices additionally \( a \leq b \iff a \wedge b = a \), 
        the rest is straightforward. 
    \end{proof}

    \begin{definition}
        A lattice is complete if each subset \( S \subset L \) has a greatest lower bound and least upper bound in \( L \),
        given by \( \bigvee_{s \in S} s \), \( \bigwedge_{s \in S} s \) respectively.
    \end{definition}
    It might appear that when a lattice is bounded, it is also complete. The two definitions coincide
    for finite \( L \), the meet and join of a lattice are only defined on pairs of elements and with that on finite \( S \), 
    while in general \( S \) might be an uncountable set. Similar distinctions should be noted for distributivity, where the counterpart 
    is completely distributive.

    \begin{definition}
        A complete distributive lattice \( L \) is join infinite distributive (JID) if it holds for any \( S \subset L \) and \( a \in L \)
        \[
            a \wedge \bigvee_{b \in S} b = \bigvee_{b \in S} a \wedge b,
        \]
        analogously \( L \) is meet infinite distributive (MID) if it holds for such \( S \) and \( a \):
        \[
            a \vee \bigwedge_{b \in S} b = \bigwedge_{b \in S} a \vee b.
        \]
    \end{definition}

    While the random variables \( \mathcal{X} : \Omega \to \mathbb{R} \) formed a ring, risk measures importantly do not with respect to 
    the same operations, e.g. \( \rho_1(X + m) + \rho_2(X + m) = \rho_1(X) + \rho_2(X) - 2m \), so cash-invariance is not maintained. 
    If risk measures are supposed to have a relation to random variables like \( \mathbb{K}[\mathbf{x}] \) has to \( \aff{\mathbb{K}}{n} \),
    as in Example \ref{algstructuresheaf}, we need a different structure, and lattices are a promising candidate.
    \begin{theorem} \label{monrmarelattice}
        The space of monetary risk measures \( \mathcal{R} \), that is functionals \( \rho : \mathcal{X} \to \mathbb{R} \), which satisfy
        the monotonicity and cash invariance from Definition \ref{defriskmeasure}, are a distributive unbounded lattice under the operations 
        \begin{align*}
            (\rho_1 \vee \rho_2)(X) &:= \max \{\rho_1(X), \rho_2(X) \} \\
            (\rho_1 \wedge \rho_2)(X) &:= \min \{\rho_1(X), \rho_2(X) \}.
        \end{align*}
        By modifying the codomain of risk measures to the extended real numbers \( \extreals = \mathbb{R} \cup \{ -\infty, \infty \} \)
        and adding the risk measures \( \rho_{\bot}(X) = -\infty,\ \rho_{\top}(X) = \infty \) to \( \mathcal{R} \) one obtains a distributive,
        bounded, complete lattice \( \extrm \) of which \( \mathcal{R} \) is a sublattice.
    \end{theorem}
    \begin{remark} \label{completionofriskmeasures}
        There is a nice, intuitive interpretation of these measures, but it comes at the cost of Lipschitz continuity, which is not too much of a burden, 
        since we have to keep only two special cases in mind. Otherwise, we do not allow \( \rho \in \extrm \) to take on the value \( - \infty \).
    \end{remark}

    \begin{lemma} \label{extrmisjidandmid}
       The lattice \( \extrm \) is JID and MID.
    \end{lemma}
    \begin{proof}
        \( \extrm \) is complete by Lemma \ref{extrmiscomplete} and the equality can be verified pointwise.
        We begin with JID. Let \( I \subseteq \extrm \), \( \rho \in \extrm \), \( X \in \mathcal{X} \)
        and let \( p = \rho(X) \) and \( q = \bigvee_{\nu \in I} \nu(X) = \bigvee_{\nu \in I} q_\nu = \sup_{\nu \in I} q_\nu \)
        be distinct members of \( \extreals \), since otherwise the result is trivial. Keep in mind that \( \extreals \) is totally ordered.
        Proceed by cases:
        \begin{itemize}
            \item \( p \geq q \): By assumption \( p \wedge q = q \) further for each \( \nu \in I \) we also have \( q_{\nu} \wedge p = q_{\nu} \)
                as otherwise we would have a contradiction. Therefore \( \bigvee_{\nu \in I} q_\nu \wedge p = \bigvee_{\nu \in I} q_\nu = q\).
            \item \( p < q \): By assumption \( p \wedge q = p \), suppose that there's no \( \nu \in I \) s.t. \( q_{\nu} > p \),
                then \( p \) is an upper bound for the set \( \{ q_{\nu} \}_{\nu \in I} \) and we would have \( \sup_{\nu \in I} q_{\nu} \leq p \)
                contradicting our hypothesis. Therefore we have some \( \nu^\ast \in I \) with \( q_{\nu^\ast} > p \) and \( \bigvee q_{\nu} \wedge p \geq q_{\nu^\ast} > p \).
        \end{itemize}
        The argument then goes analogously for MID.
    \end{proof}

    To give a taste of the intuitiveness of our construction and why lattices are a natural analog for algebraic geometry, let us turn to the geometric side:
    \begin{corollary} \label{acceptanceoperations}
        The operations in \( (\extrm, \vee, \wedge) \) correspond to intersections and unions of acceptance sets
        and vice-versa. Further the lattice order translates into acceptance sets as: \( \rho_1 \leq \rho_2 \) for \( \rho_1, \rho_2 \in \extrm \) if and only if \( \mathcal{A}_{\rho_2} \subseteq \mathcal{A}_{\rho_1} \).
    \end{corollary}
    \iffalse
    As C.G. Jung already said, "Life calls not for perfection, but for completeness" we shall 
    focus our study on \( \extrm \) from hereon as complete lattices are better behaved 
    regarding algebraic analysis. A hint at this is Definition \ref{atomisticlattice}, and as we will see, the definition or verification of ideals becomes easier if we grant ourselves this freedom. 
    \fi
\end{subsection}

\begin{subsection}{Ideals and Filters in Lattices} 
The subject of this section is the analog of ideals and their properties, such as being prime or principal in rings for distributive lattices. 
There are actually two structures in lattices that fit the concept: ideals and filters, and both will be relevant.
\begin{definition}
        An ideal of a semilattice \( (L, \vee) \) is a non-empty set \( I \subseteq L \) s.t.
        \begin{itemize}
            \item \( a, b \in I \implies a \vee b \in I \) 
            \item \( a \in I \) and \( b \leq a \implies b \in I \).
        \end{itemize}
        If \( L \) is bounded with element \( e \), we additionally require \( e \in I \).
        The dual notion of an ideal of a bounded semilattice is a filter \(F \subseteq L \) with \( F \neq \emptyset \), namely
        \begin{itemize}
            \item \( a, b \in F \implies a \wedge b \in F \) 
            \item \( a \in F \) and \( b \geq a \) imply \( b \in F \).
        \end{itemize}
        If \( L \) is bounded, we also require \( e \in F \).
        An ultrafilter \( U \) is a filter, that is proper and is not contained in any other proper filter.
        An ideal \( P \subset L \) is a prime ideal if its complement is a filter.
        For a bounded lattice, a set \( F \subset L \) is a filter if it is one for \( (L, \vee, \bot) \),
        similarly an ideal of a bounded lattice is one for \( (L, \wedge, \top) \).
        In particular, the second defining property for ideals is equivalent 
        to the more familiar \( a \in I, b \in L \implies a \wedge b \in I \).
    \end{definition}
    \begin{remark}
        For further use, denote by \( \text{Spec } L \) the set of prime ideals of a semilattice, also called the spectrum of \( L \).
    \end{remark}

    \begin{proposition} \label{propositionprime}
        For \( L \) a distributive lattice an ideal \( P \subset L \) being prime is equivalent to:
        \begin{itemize}
            \item  \( P \) is proper (\( P \subsetneq L \)) and \( a \wedge b \in P \implies a \in P \text{ or } b \in P \).
        \end{itemize}
        If \( L \) is also bounded, \( P \) being prime is equivalent to the following: 
        \begin{itemize} 
            \item \( \top \notin P \) and \( a \wedge b \in P \implies a \in P \text{ or } b \in P\).
            \item \( P \) is the kernel of a lattice homomorphism \( f : A \to \{0, 1\} \), where \( \{0, 1\} \) is the binary lattice.
        \end{itemize}
    \end{proposition}
    \begin{proof}
        Theorem 9.8 of \cite{birkhoff1948lattice}.
    \end{proof}
     Especially the first property reminds us of classical prime ideals from rings. 
        Standard examples of an ideals are sets of the form \( (-\infty, a] \) in \( \mathbb{R} \), and the reader may convince themself that it is also prime (yielding a specimen of a filter too). 
        Filters are a very flexible tool that can be used to give an alternative axiomatic foundation for topology by leveraging the fact that the power set forms a lattice with intersections as the meet and union as the join.
        Ultrafilters are the dual concept of maximal ideals and encode in the topology setting convergence to a point.

    \begin{proposition} \label{meetirreducibleprime}
        Let \( L \) be a distributive lattice, then the principal ideal \( (a)^\downarrow  \)
        is prime if and only if \( a \) is meet irreducible,
        meaning there are no \( b, c \in L \setminus \{a\} \) s.t. \( a = b \wedge c \).
    \end{proposition}
    \begin{proof}
        Exercise 9.6.4(a) in \cite{birkhoff1948lattice}.
    \end{proof}
    We will later study "meet irreducible" elements of \( \extrm \), although the reader is invited to already make their own guesses.
    There will be a natural interpretation, what makes them irreducible in terms of risk.

    \begin{proposition} \label{separationlatspec}
        In a distributive lattice \( L \), every ideal is contained in a prime ideal.
        Moreover prime ideals can distinguish points of the lattice, meaning for distinct \( a, b \in L \),
        there is a prime ideal \( P \) s.t. only one of \( a,b \) is in \( P \).
    \end{proposition}
    \begin{proof}
        Corollary 2.1.17 in \cite{genlattheorygraetzer}.
    \end{proof}
    \begin{remark}
        A consequence of Proposition \ref{separationlatspec} is that if we endow 
        \( \spec{L} \) with the direct analog of the Zariski topology,
        meaning a subbasis of open sets is given by \( D(a) := \{ \mathfrak{p} \in \spec{L} | a \notin \mathfrak{p} \} \)
        for \( a \in L \), referred to as Stone topology, then \( \spec{L} \) is a \( T_0 \)-space.
        In classical lattice theory, this is the basis for a series of celebrated representation theorems
        such as that of Stone and Birkhoff. We won't leverage this topology on \( \spec \extrm \), because
        there are far more prime ideals than our use cases would permit.
    \end{remark}
\end{subsection}

%\begin{subsection}{Quotient Lattices} \label{quotientlatticesect}
    In the algebro-geometric setting, we defined coordinate rings for a variety  \( V \)
    via \( \mathbb{K}[X]/J \), where \( J = I(V) \) was the vanishing ideal.
    Moreover, for each Zariski open subset \( U \subset V \) we associated
    the ring of regular functions and at a specific point \( p \in V \), the 
    local ring \( \mathcal{O}_{p} \) and the residue field \( \resf{p} \).
    The natural question is if there are similar concepts for the lattice \( \mathcal{R} \)
    and acceptance sets of members. As it turns out there is a less unique choice 
    in this setting, because both ideals and filters allow for similar constructions. 
    We will cover both here and later state in Remark \ref{dualfilterquotient}, how the constructions 
    hold for both of them.

    Let \( I \) be an ideal in a distributive lattice \( (L, \vee, \wedge) \),
    and define the relation \( a \sim b \iff \exists i \in I,\ a \vee i = b \vee i\).
    Naturally for any \( a \in L \) and \( i \in I \), \( a \vee i = a \vee i \),
    so \( (\cdot \sim \cdot) \) is reflexive. By commutativity of \( \vee \) the relation is also symmetric
    and if \( a \sim b, b \sim c \) with corresponding \( i_a, i_c \) one derives
    \begin{align*}
        a \vee i_a &= b \vee i_a \\
        i_c \vee (a \vee i_a) &= i_c \vee (b \vee i_a) \\
        a \vee (i_a \vee i_c) &= c \vee (i_c \vee i_a)
    \end{align*}
    by associativity and commutativity of \( \wedge \), hence \( a \sim c \) with \( i_c \vee i_a \)
    which is in \( I \) as \( I \) is an ideal.
    We conclude that \( (\cdot \sim \cdot) \) is an equivalence relation.
    Moreover, let \( a_1 \sim b_1 \), \( a_2 \sim b_2 \) via \( i_1, i_2 \in I \)
    then 
    \begin{align*}
        (i_1 \vee a_1) \vee (i_2 \vee a_2) = (i_1 \vee b_1) \vee (i_2 \vee b_2) \\
        (a_1 \vee a_2) \vee (i_1 \vee i_2) = (b_1 \vee b_2) \vee (i_1 \vee i_2)
    \end{align*}
    so \( (a_1 \vee a_2) \sim (b_1 \vee b_2) \). Similarly
    \begin{align*}
        (i_1 \vee i_2) \vee \left( (i_1 \vee a_1) \wedge (i_2 \vee a_2) \right)  \\
     \left ( (i_1 \vee i_2) \vee (i_1 \vee a_1) \right) \wedge \left( (i_1 \vee i_2) \vee (i_2 \vee a_2) \right)  \\
     \left ( (i_1 \vee i_2) \vee a_1) \right) \wedge \left( (i_1 \vee i_2) \vee a_2) \right) \\
     (i_1 \vee i_2) \vee (a_1 \wedge a_2)
    \end{align*}
    and doing the same with \( (i_1 \vee b_1) \wedge (i_2 \vee b_2) = (i_1 \vee a_1) \wedge (i_2 \vee a_2) \), gives 
    with \( i_1 \vee i_2 \in I  \) (since it is an ideal) that \( (a_1 \wedge a_2) \sim (b_1 \wedge b_2) \), making \( (\cdot \sim \cdot) \) 
    a compatible relation with the lattice operations and by the quotient theorem for universal algebras,
    we have a well defined "quotient lattice" \( L / I \). On a side note, if \( \vee, \wedge \) are preserved
    so are any inequalities between elements in the lattice under the equivalence relation.
    \begin{definition} \label{defquotlat}
        For an ideal \( I \) of a distributive lattice \( (L, \vee, \wedge) \)
        we call the lattice \( (L / I, \vee, \wedge) \) the quotient lattice of \( L \)
        by \( I \) (or just quotient lattice if it is clear from context).
    \end{definition}

    What does that mean in the case of risk measures?
    By Proposition \ref{acceptancesetareideals} any acceptance set induces an ideal \( \acceptidealb{A} \).
    Clearly all elements in \( \acceptidealb{A} \) are equivalent to each other in \( \extrm / \acceptidealb{A} \),
    and since \( \rho_\bot \in \acceptidealb{A} \) it is sensible to summarize them under this representant \( [ \rho_\bot ] \).
    All elements \( \rho \in \extrm \) which dominate members of \( \eta \in \acceptidealb{A} \) in the sense \( \eta \leq \rho \)
    satisfy \( \rho \vee \eta = \rho \), hence they are in distinct classes.
    The most interesting phenomenon happens for those \( \rho \in \extrm \), which neither satisfy \( \rho \leq \eta \)
    nor \( \rho \geq \eta \). For example if two risk measures \( \rho_1, \rho_2 \) are different on some subset \( S \subset \mathcal{X} \)
    but otherwise coincide, they could become equivalent if there's some \( \eta \in \acceptidealb{A} \), which dominates both of them on \( S \).
    Additionally since they are not in the ideal, we know that there are \( X, Y \in A \), s.t. \( \rho_1(X) > 0, \rho_2(Y) > 0 \),
    but \( \eta(X) \leq 0, \eta(Y) \leq 0 \), so if \( \rho_1 \sim \rho_2 \) then they must coincide for all \( X \in A \) which they
    reject, thus we could still evaluate equivalence classes of risk measures \( \rho \) with respect to their rejection sets.
    %A similar observation holds also for \( \extrm / \rejectidb{A} \) and the equi
    \newline
    For concrete computations this is somewhat inadequate, since we analyze risk measures both with respect to their values they accept and those
    they don't. In general just because two risk measures \( \rho_1, \rho_2 \in \extrm \) agree on their rejection set when restricted \( A \), 
    does not imply a globally the existence of some \( \eta \in \acceptidealb{A} \) with \( \eta \vee \rho_1 = \eta \vee \rho_2 \). 
    The most valuable information which therefore can still be extracted from this setting is given less by concrete values and more by families
    of risk measures given either by ideals and filters.
    \newline
    \newline
    The behavior of ideals from our larger lattice \( \extrm \) should then translate into ideals of \( \extrm / \acceptidealb{A} \), similarly
    how ideals from a general ring translate into ideals of a quotient ring with respect to some ideal.
    Also we would like to have a relation between \( \text{Spec } \extrm \) and \( \text{Spec } \extrm / \acceptidealb{A}\).
    Fortunately, this is a pattern which is shared by distributive lattices. 
    \begin{definition}
        Let \( L \) be a lattice and denote by \( \idl{L} \) the set of ideals in \( L \)
        and by \( \flt{F} \) the set of filters of \( L \).
    \end{definition}
    \begin{proposition}
        \( \idl{L} \) is itself a distributive lattice under the operations
        \begin{align*}
            I \wedge J = \{ i \wedge j | i \in I, j \in J \} \\
            I \vee J = \{ i \vee j | i \in I, j \in J \}
        \end{align*}
        for \( I, J \in \idl{L} \). Moreover, the meet simplifies to setwise intersection
        \begin{align*}
            I \cap J = I \wedge J
        \end{align*}
        and if \( L \) was bounded then \( \idl{L} \) is complete.
    \end{proposition}
    \begin{proof}
        See Lemma 2.5.1, Corollary 1.3.2, Corollary 1.3.15 in \cite{genlattheorygraetzer}.
    \end{proof}
    \begin{lemma} \label{lemmaprincipalequivalence}
        Let \( (i)^\downarrow \subset L \) be a principal ideal with generator \( i \) in a distributive lattice \( L \).
        Then for \( a, b \in L \) it holds
        \[
            [a] = [b] \iff a \vee i = b \vee i
        \]
    \end{lemma}
    \begin{proof}
        Assuming the RHS implies the LHS by Definition. For the converse, let \( j \in (i)^\downarrow \),
        s.t. \( a \vee j = b \vee j \), applying the join with \( i \) to both sides then reduces to the statement.
    \end{proof}
    \begin{proposition} \label{correspondencetheorem}
        For distributive lattices the correspondence theorem holds, that is if \( I \in \idl{L} \)
        and \( (I)^\uparrow = \{ J \in \idl{L} | I \subseteq J \} \), then we have a bijection
        \[
            \Phi : (I)^\uparrow \mapsto \idl{L/I},
        \]
        which preserves inclusion, primeness and maximality.
    \end{proposition}
    \begin{proof}
        Being bijective is a restatement of Lemma 2.6.11 in \cite{genlattheorygraetzer} (second isomorphism theorem for universal algebras) by noting that one can substitute 
        for congruence relation the equivalence relation defined by ideals for their quotient lattices.
        As inclusion is preserved under maps, it is in particular by the projection \( \phi : L \to L / I \)
        and then \( \Psi \) being bijective maintains the inclusion of the ideals.
        If inclusions are preserved so is the property of being maximal.
        Further it holds \( \pi(J) = \pi(J)^\downarrow \) (the generated ideal by \(J\) in the quotient), because 
        if \( [a] \leq [j] \) for some \( j \in \pi(J) \), we can write it as \( [a] = [j \wedge b] \) for some \( b \in L \)
        but \( J \) is an ideal so \( j \wedge b \in J \) and then we have \( \pi^{-1}([a]) \in J \).
        Let \( P \in (I)^\uparrow \) be prime, then if \( [a \wedge b] \in \pi(P) \) also \( a \wedge b \in \pi^{-1}([a \wedge b]) \)
        is also in \( P \) and by assumption then \( a \) or \( b \) too.
    \end{proof}
    \begin{proposition} \label{quotientprojectionlattice}
        For a distributive lattice \( L \) and \( I, J \in \idl{L} \) s.t. \( I \subseteq J \), there is a well-defined lattice homomorphism
        \[
            \varphi_{J,I} : L / I  \mapsto L / J, \quad [x]_{I} \mapsto [x]_{J},
        \]
        for \( x \in L \).
    \end{proposition}
    \begin{proof}
        Let \( a,b \in L \) and \( [a]_{I}, [b]_{I} \in L / I \) be their projections into \( L / I \) and assume that \( [a]_{I} = [b]_{I} \),
        by Definition there's some \( i \in I \) s.t. \( a \vee i = b \vee i \), since \( I \subseteq J \), we also have that \( [a]_{J} = [b]_{J} \)
        making \( \varphi \) well-defined as a map. That it is also a homomorphism is a standard verification.
    \end{proof}
    Essentially, if \( I \subseteq J \) the equivalence relation of \( J \) is broader than the one of \( I \). Note that 
    \( \acceptidealb{U} \) was inclusion reversing (Corollary \ref{acceptidealincrev}) , so for \( U \subseteq V \) 
    the ideals satisfy \( \acceptidealb{U} \subseteq \acceptidealb{V} \). The intuition is that accepting \( V \) is less 
    restrictive then it is to accept \( U \), so the ideal is larger. This motivates the next observation, which should remind the reader 
    of something.
    \begin{lemma} \label{triangleidealcommutes}
        Let \( I \subseteq J \subseteq K \) be ideals in a distributive lattice \( L \), then the following 
        diagram commutes:
        \[
        \begin{tikzcd}[row sep=large, column sep=large]
            L / I \arrow[dr, "\varphi_{J, I}"] \arrow[rr, "\varphi_{K, I}"] & & L / K \\
                                                                            & L / J \arrow[ur, "\varphi_{K, J}"]  
        \end{tikzcd}
        \]
    \end{lemma}
    \begin{proof}
        By Proposition \ref{quotientprojectionlattice} the projections are all well defined
        and we need to show for \( x \in L \) with \( \varphi_{I}(x) = [x]_I \), 
        and \( \varphi_{J, I}([x]_{I}) = [x]_{J} \), that \( \varphi_{K, J}([x]_{J}) = \varphi_{K, I}([x]_{I}) \).
        But this holds by Definition, since \( \varphi_{K, J}([x]_{J}) = [x]_{K} \) and \( \varphi_{K, I}([x]_{I}) = [x]_{K} \).
    \end{proof}
    While ideals have a straightfoward relation, which is symmetric to the situation for rings and their quotients, the case of  
    filters in \( L / I \) requires more care.
    For this purpose the following lemma comes in handy:
    \begin{lemma} \label{filterandidealsdonotcontain}
        Let \( L \) be a distributive lattice, for any \( I \in \idl{L} \)
        and \( F \in \flt{L} \), neither \( I \subseteq F \) nor \( F \subseteq I \) holds.
    \end{lemma}
    \begin{proof}
        Suppose the contrary, i.e. \( I \subseteq F \), by Definition
        any \( I \) is downward closed and absorbing with respect to \( \wedge \), while \( F \)
        is upward closed and absorbing with respect to \( \vee \). Take any \( a \in L \)
        and \( i \in I \), if \( a \) is not already in one of them
        then \( a \wedge i \in I \subseteq F \), but \( a \wedge i \leq a \) and \( F \)
        is upward closed, hence \( a \in F \) and by arbitrariness of \( a \), \( F \) can not be proper
        and therefore not a filter.
        The converse for \( F \subseteq \) follows dually with \( a \vee f \) for \( f \in F \).
    \end{proof}
    Note that this does not necessarily mean that all filters and ideals are disjoint! Those who do however
    have a special relation.
    \begin{proposition} \label{filtersinidealquot}
        Let \( I \in \idl{L} \) then there's an inclusion preserving bijection 
        \[ \Psi : \satflt{I}{L} \to \flt{L / I}, \]
        where \( \satflt{I}{L} \) is the set of filters, which are
        disjoint and saturated with respect to I.
        A filter \( F \) is saturated with respect to an ideal \( I \), if
        for \( a \in F \) and \( b \in L \), 
        \[ 
            [a]_{L / I} = [b]_{L / I} \implies b \in F.
        \]
        For such filters, primeness and being an ultrafilter is preserved.
    \end{proposition}
    \begin{proof}
        Let \( \varphi : L \to L / I \) be the canonical projection and define \( \Psi(F) := \varphi(F) \) in the sense of sets.
        \begin{itemize}
            \item \( F \in \satflt{I}{L} \implies \Psi(F) \in \flt{L / I} \).
                \begin{itemize}
                    \item 
                        Take \( F \in \satflt{I}{L} \) with \( a, b \in F \)
                        then \( \varphi(a) \wedge \varphi(b) = \varphi(a \wedge b) \)
                        so for \( [a], [b] \in \Psi(F) \) also \( [a] \wedge [b] = [a \wedge b] \in \Psi(F) \).
                        Similarly for \( a \in F \) and \( b \in L \), we can argue \( [a \vee b] \in \Psi(F) \).
                        Further for \( a \in F \) and \( b \in L \), \( [a] \leq [b] \iff [a \vee b] = [b] \)
                        and \( F \) is a filter, so \( a \vee b \in F \), but then by saturation \( b \in F \), so \( [b] \in \Psi(F) \).
                \end{itemize}
            \item \( F \in \satflt{I}{L} \) is prime, then \( \Psi(F) \in \flt{L / I} \) is as well.
                \begin{itemize}
                    \item 
                        By the previous \( \Psi(F) \in \flt{L / I} \).
                        If \([a \vee b] \in \Psi(F) \) for \( a, b \in L \) then there are \( f \in F \) and \( i \in I \)
                        s.t. \( (a \vee b) \vee i = f \vee i \), \( F \) is absorbing with respect to \( \vee \), so \( f \vee i \in F \)
                        and therefore \( (a \vee b) \vee i \in F \). Now \( F \) is prime and disjoint from \( I \), so only \( (a \vee b) \in F \)
                        is possible and applying primeness again gives \( a \in F \) or \( b \in F \), but then one of their equivalence classes is in \( \Psi(F) \).
                \end{itemize}
            \item \( F' \in \flt{L / I} \implies \Psi^{-1}(F') \in \flt{L} \) and saturated with \( I \).
                \begin{itemize}
                    \item 
                        Let \( F' \in \flt{L / I} \), if \( a, b \in \Psi^{-1}(F') \) then \( [a], [b] \in F' \)
                        and as \( F' \) is a filter, \( [a \wedge b] \in F' \) but then \( a \wedge b \in \Psi^{-1}(F') \).
                        Similarly for \( a \in \Psi^{-1}(F') \) and \( b \in L \), it holds \( [a \vee b] \in F' \),
                        and we conclude \( a \vee b \in \Psi^{-1}(F') \). Combining the results we have that \( \Psi^{-1}(F') \) is a filter.
                        Let \( a \in \Psi^{-1}(F') \) and \( b \in L \), s.t. 
                        \( [a] = [b] \in F' \), trivially then also \( b \in \Psi^{-1}(F') \), making \( \Psi^{-1}(F') \) saturated.
                \end{itemize}
            \item \( F' \in \flt{L / I} \) is prime then \( \Psi^{-1}(F') \in \flt{L} \) is as well.
                \begin{itemize}
                    \item 
                        Suppose that \( F' \in \flt{L / I} \) is also prime, by the previous \( \Psi^{-1}(F') \in \satflt{I}{L}\),
                        assume \( a \vee b \in \Psi^{-1}(F')\), then \( [a \vee b] \in F' \) and due to assumption \( [a], [b] \in F' \)
                        by saturation of \( \Psi^{-1}(F') \) it follows that \( a \in \Psi^{-1}(F') \) or \( b \in \Psi^{-1}(F') \).
                \end{itemize}
            \item \( F' \in \flt{L / I} \implies \Psi^{-1}(F') \in \satflt{I}{L} \).
                \begin{itemize}
                    \item 
                        Assume for contradiction that there is an \( i \in I \cap \Psi^{-1}(F') \), then \( [i] \in F' \) and as
                        all elements in \( I \) are in the same quotient class, also \( [j] \in F' \) for all \( j \in I \setminus \{i\}\),
                        but then \( I \subseteq \Psi^{-1}(F') \), and \( \Psi^{-1}(F') \) can not have been a filter according to Lemma \ref{filterandidealsdonotcontain}
                        and therefore \( I \) and \( \Psi^{-1}(F') \) are disjoint.
                \end{itemize}
        Hence \( \Psi \) is well defined bijective map with the asserted domain and codomain.
        In the proof we defined \( \Psi \) even on the level of sets instead of filters hence it
        is inclusion preserving and therefore ultrafilters in \( \satflt{I}{L} \) are preserved.
        \end{itemize}
    \end{proof}
    \begin{remark} \label{dualfilterquotient}
        All of the previous constructions also exist for filters, which one might argue is an immediate consequence by
        the duality principle (see also the end of Section 1.1 of \cite{genlattheorygraetzer}). 
        If \( F \subset L \) is a filter, then 
        we can construct the quotient lattice \( L / F \) under the equivalence relation 
        \[
            \rho_1 \sim \rho_2 \iff \exists \eta \in F, \rho_1 \wedge \eta = \rho_2 \wedge \eta.
        \]
        Then everything in \( F \) is sent to \( \top \) in \( L / F \). Additionally all filters
        containing \( F \) are sent to filters in \( L / F \) and the properties of being an ultrafilter
        or prime (that is \( a \vee b \in F \iff a \in F \text{ or } b \in F \)) are maintained. 
        Also for filters \( F_1 \subseteq F_2 \) there is a well-defined lattice homomorphism \( \varphi : L / F_2 \to L / F_1 \)
        and disjoint saturated ideals from a filter correspond to ideals in the filter quotient.
    \end{remark}
\end{subsection}

%\begin{subsection}{A Sheaf on Monetary Risk Measures}
    Let \( \mathcal{X} = L^\infty(\Omega, \Sigma, \mathbb{P}) \) with its usual topology for this section.
    Define a contravariant functor based on the map \( \rejectidb{\cdot} \) from Definition \ref{rejectionsetsarefilters}
    with the signature \( \presh :\textbf{Open}(\mathcal{X}) \to \quotlat{\extrm}, U \mapsto \frac{\extrm}{\rejectidb{U}} \), where we understand 
    \( \textbf{Open}(\mathcal{X}) \) as the category with objects the open sets \( U \subseteq \mathcal{X} \) and 
    and morphisms being inclusions \( U_1 \subseteq U_2 \) of them. The codomain is the category with objects being the quotient lattices of \( \extrm \)
    and morphism induced as in Remark \ref{dualfilterquotient}.
    \begin{theorem} \label{rmpresheaf}
        The contravariant functor \( \presh : \textbf{Open}(\mathcal{X}) \to \quotlat{\extrm} \) defines a separated presheaf on
        \( \mathcal{X} \).
    \end{theorem}
    \begin{proof}
        \begin{itemize}
            \item \( \presh \) is a presheaf.
                \begin{itemize}
                    \item
                        If \( U = \emptyset \) then \( \rejectidb{U} = \extrm \) and \( \frac{\extrm}{\rejectidb{U}} = \{0\} \) the trivial lattice, clearly all lattices have a lattice morphism 
                        to \( \{0\} \).
                        For each open \( U \subseteq \mathcal{X} \) with the identity morphism of sets there's also an identity morphism of lattices \( \frac{\extrm}{\rejectidb{U}} \).
                        For open \( V \subseteq U \) we get by Lemma \ref{rejectionsetsarefilters}, that \( \rejectidb{U} \subseteq  \rejectidb{V} \) as filters in \( \extrm \)
                        and Remark \ref{dualfilterquotient} then gives a morphism \( \varphi_{V,U} : \frac{\extrm}{\rejectidb{U}} \to \frac{\extrm}{\rejectidb{V}}, [\rho]_U \mapsto [\rho]_V \),
                        by Lemma \ref{triangleidealcommutes} it also holds for open \( W \subseteq V \) that \( \varphi_{W,U} = \varphi_{W,V} \circ \varphi_{V,U} \).
                \end{itemize}
            \item \( \presh \) is separated.
                \begin{itemize}
                    \item
                        Take any \( [\rho_1], [\rho_2] \in \presh(U) \), where \( \rho_1, \rho_2 \in \rejectidb{U} \), with \( [\rho_1]|_{U_i} = [\rho_2]|_{U_i} \)
                        for all \( U_i \). This implies for each \( i \) some \( \eta_i \in \rejectidb{U_i} \) s.t. \( \rho_1 \wedge \eta_i = \rho_2 \wedge \eta_i \). 
                        Note that \( \rho_1, \rho_2, \eta_i \in \extrm \), so we can define the risk measure \( \eta = \bigvee_{i} \eta_i \). Since \( \eta \geq \eta_i \) 
                        for all \( i \), it follows \( \eta(X) > 0 \) for all \( X \in U \), so \( \eta \in \rejectidb{U} \). Moreover,
                        \( \rho_1 \wedge \eta = \bigvee_{i} (\rho_1 \wedge \eta_i) = \bigvee_{i} (\rho_2 \wedge \eta_i) = \rho_2 \wedge \eta \), where we used Lemma \ref{extrmisjidandmid}
                        for the equalities, but then \( [\rho_1] = [\rho_2] \).
                \end{itemize}
            
        \end{itemize}
    \end{proof}
    While we were previously working mostly with \( \acceptidealb{U} \), it is quite natural to use \( \rejectidb{U} \) for the presheaf instead.
    First of all acceptance sets are often not open, but closed, while "rejection sets" are typically open. Also we get now a close analogy
    to the local rings. We discovered that every \( X \in \mathcal{X} \) defines an element \( \acceptidealb{\{X\}} \in \text{Spec } \extrm \), but at the same time 
    they also define a filter via \( \rejectidb{\{X\}} \in \text{Spec } \extrm \) and it even holds \( \extrm \setminus \rejectidb{\{X\}} = \acceptidealb{\{X\}} \).
    \begin{proposition}
        \( \acceptidealb{\{X\}} \) is under the canonical map, \( \varphi : \extrm \to \frac{\extrm}{\rejectidb{\{X\}}} \) mapped to the unique maximal ideal
        of \( \frac{\extrm}{\rejectidb{\{X\}}} \).
    \end{proposition}
    \begin{proof}
        By Remark \ref{dualfilterquotient}, \( \varphi(\rejectidb{\{X\}}) = [\top] \), and as \( \acceptidealb{\{X\}} \) is the complement of \( \rejectidb{\{X\}} \)
        all remaining elements of \( \frac{\extrm}{\rejectidb{\{X\}}} \) are in its image.
    \end{proof}
    So even though we did not have maximal ideals for each point \( X \in \mathcal{X} \) in \( \text{Spec } \extrm \), they still define unique maximal ideals in the quotient lattices
    with respect to their filters, as do maximal ideals \( \mathfrak{p} \in \text{Spec } A \) for a commutative ring \( A \) in their local rings \( A_{\mathfrak{p}} \).
    Therefore we could interpret such filters as analogues to the multplicative sets \( A \setminus \mathfrak{p} \) and were granted all that despite not having a concept of inverses in \( \extrm \).
    \newline
    Also we can now answer what the stalks of our presheaf are:
    \begin{corollary}
        The stalk \( \presh_X \) at a point \( X \in \mathcal{X} \) of a presheaf \( \presh \) on the topological space \( \mathcal{X} \),
        is defined as the colimit of the image of the open neighborhoods of the point:
        \[ 
           \presh_{X} := \varinjlim_{x \in U} \presh(U),
        \]
        For the presheaf from Theorem \ref{rmpresheaf} the stalk satisfies
        \[
            \presh_{X} = \frac{\extrm}{\rejectidb{\{X\}}}.
        \]
    \end{corollary}
    \begin{proof}
        By Proposition \ref{colimitforquotientlattice} it suffices to verify 
        that \( \rejectidb{\{X\}} = \cup_{x \in U} \rejectidb{U} \).
        Take \( \rho \in \cup_{x \in U} \rejectidb{U} \), because \( X \in U \), we also have \( \rho(X) > 0 \)
        and \( \rho \in \rejectidb{\{X\}} \).
        Conversely, let \( \rho \in \rejectidb{\{X\}} \), then \( \rho \neq \rho_{\bot} \) and if \( \rho = \rho_{\top} \)
        then it is already contained in all \( \rejectidb{U} \), so we can focus on \( \rho \in \mathcal{R} \). 
        By Definition \( c := \rho(X) > 0 \), since \( \rho \in \mathcal{R} \) is Lipschitz continuous, there exists
        \( \delta > 0 \), s.t. \( \norm{Y - X}_{\infty} < \delta \implies |\rho(Y) - \rho(X)| < \frac{c}{2} \),
        note that if \( \rho(Y) \leq 0 \) then from \( c = \rho(X) > 0 \), the absolute value can not be smaller than \( \frac{c}{2} \)
        (in fact no smaller than \( c \)), so \( \rho(Y) > 0 \) for all \( Y \in B_{\delta}(X)\).
        But this means that \( \rho \) is strictly positive on some open neighborhood of \( X \), hence \( \rho \in \cup_{x \in U} \rejectidb{U} \).
    \end{proof}
    Let \( \textbf{BDLat} \) denote the category of bounded distributive lattices.
    \begin{proposition} \label{sheafification}
        %: \text{im } \presh \to \Pi \presh_x
        There is a monomorphism \( \text{sh} : \presh \to \mathcal{F} \) of presheaves, from the presheaf \( \presh \) of Theorem \ref{rmpresheaf}
        to a sheaf \( \mathcal{F} : \textbf{Open}(\mathcal{X}) \to \textbf{BDLat} \) on \( \mathcal{X} \), given by
        \[
            \mathcal{F}(U) \mapsto \left \{ (\sigma_X)_{X \in U} \in \Pi_{X \in U} \presh_X | (\sigma_X)_{X \in U} \text{ satisfies } (*)  \right \},
        \]
        where \( (*) \) is the property that for every \( X \in U \) there exists some open neighborhood \(V \subseteq U\) of \( X \) and a section \( [\rho] \in \presh(V) \)
        s.t. for all \( Y \in V \), it holds \( \sigma_Y = \varphi_{\{Y\},V}([\rho]) \).
    \end{proposition}
    \begin{proof}
        %Use Lemma 6.17.5 of \cite{stackprojectsheafification} and 
        The categorical product of distributive bounded lattices coincides with the Cartesian product \cite{catprodlattice}
        of them, so \( \prod_{X \in U} \presh_X \) is the lattice where \( \wedge, \vee \) are defined componentwise on tuples
        \( ([\rho]_{\presh_X})_{X \in U} \). By the \( (\ast) \) property of \( (\sigma_X)_{X \in U}, (\delta_X)_{X \in U} \in \mathcal{F}(U) \) 
        we have for any \( X \in U \) open neighborhoods \( V_{\sigma}, V_{\delta} \) over which for each \( Y \in V_{\sigma}, Z \in V_{\delta} \) it holds 
        \begin{align*}
            \sigma_Y = \varphi_{\{Y\},V_{\sigma}}([\rho]) = [\rho]_{\rejectidb{Y}} \\
            \delta_Z = \varphi_{\{Z\},V_{\delta}}([\eta]) = [\eta]_{\rejectidb{Z}}
        \end{align*}
        Since this holds over the whole neighborhood we bend the notation a little and write \( \sigma|_{V_{\sigma}} = [\rho]_{\rejectidb{V_{\sigma}}} \),
        \( \delta|_{V_{\delta}} = [\eta]_{\rejectidb{V_{\delta}}} \).
        By construction \( X \in V_{\sigma} \cap V_{\delta} \), so we have an open non-empty intersection \( V' \) on which we can project 
        \( [\rho] \), \( [\eta]_{} \) and take the meet and join there. In the succeeding we express this with the meet but the argument with the join is identical.
        Therefore locally one can express the meet as:
        \begin{align} \label{eqwedgeinsheaf}
            (\sigma_X)_{X \in V'} \wedge (\delta_X)_{X \in V'} = 
            \varphi_{V', V_{\sigma}}([\rho]) \wedge \varphi_{V', V_{\sigma}}([\eta]) = [\eta \wedge \rho]_{V'}.
        \end{align}
        Let \( (\gamma_{X})_{X \in U} := (\sigma_X)_{X \in U} \wedge (\delta_X)_{X \in U}\), to verify that \( \mathcal{F}(U) \)
        is a sublattice of \( \prod_{X \in U} \presh_X \), 
        we need to show that \( (\gamma_{X})_{X \in U} \) still satisfies \( (\ast) \), but this follows from the previous, since for each \( X \in U \) 
        we constructed an open \( V' \) on which \( (\gamma_{X})_{X \in U} \) restricts for each \( Y \in V' \) to the expression from Equation \ref{eqwedgeinsheaf}.
        Any sublattice of a distributive lattice is also distributive and the bounds are of course
        \(\top := ([\top]_{\presh_X})_{X \in U}, \bot := ([\bot]_{\presh_X})_{X \in U} \).
        \newline
        \newline
        We will proceed with verifying the monomorphism but leave the sheaf property of \( \mathcal{F}(U) \) to the reader (who alternatively can convince
        themself of it by looking up the official term of this subject, provided after the proof).
        To be a morphism the following diagram should commute for \( U \subseteq \mathcal{X} \) and open \( V \subseteq U \):
        \[
        \begin{tikzcd}[row sep=large, column sep=large]
            \presh(U) \arrow[r, "\text{sh}"] \arrow[d, "\varphi_{V,U}"]& \mathcal{F}(U) \arrow[d, "\psi_{V,U}"] \\
            \presh(V) \arrow[r, "\text{sh}"] & \mathcal{F}(V),
        \end{tikzcd}
        \]
        where \( \psi_{V,U} \) is the restriction in \( \mathcal{F}(U) \), given by \( ([\rho])_{\presh_X})_{X \in U} \mapsto ([\rho])_{\presh_X})_{X \in V} \).
        %dropping all indices \( X \in U \setminus V \) of \( (\sigma_X)_{X \in U} \).
        Take \( [\rho] \in \presh(U) \) , s.t. we restrict
        \( [\rho] \) onto \( \presh(V) \), via \( \varphi_{V,U} : \extrm / \rejectidb{U} \to \extrm / \rejectidb{V} \).
        Apply then the map to get \( \text{sh} \circ \varphi_{V,U}([\rho])) = ([\rho]_{\presh_X})_{X \in V} \),
        but pointwise this is simply \( \varphi_{\{X\},V} : \extrm / \rejectidb{V} \to \extrm / \rejectidb{\{ X \} } \),
        so we have the equality 
        \[
            \varphi_{\{X\},V} \circ \varphi_{V,U}([\rho]) = (\text{sh} \circ \varphi_{V,U}([\rho])))_X,
        \]
        where the subscript \( X \) denotes the component of \( \text{sh} \circ \varphi_{V,U}([\rho]) \in \mathcal{F}(V) \) at \( X \).
        Since \( \presh \) is a presheaf the LHS simplifies:
        \[
         \varphi_{\{X\}, U}([\rho]) = (\text{sh} \circ \varphi_{V,U}([\rho])))_X
        \]
        which is precisely \( (\psi_{V,U} \circ \text{sh}([\rho]))_X \) and as \( X \) was arbitrary, the diagram commutes.
        That the morphism is injective follows, because if
        we have \( ([\rho_1]_{\presh_X})_{X \in U} = ([\rho_2]_{\presh_X})_{X \in U} \), we get for each \( X \) an \( \eta_X \in \rejectidb{\{X\}} \)
        s.t. \( \rho_1 \wedge \eta_X = \rho_2 \wedge \eta_X \) and similar to the proof from Theorem \ref{rmpresheaf} for separatedness we get 
        an \( \eta = \bigvee_{X \in U} \eta_X \) in \( \rejectidb{U} \) which satisfies \( \rho_1 \wedge \eta = \rho_2 \wedge \eta \).
    \end{proof}
    The previous map is in fact the sheafification of \( \presh \) whose concatentation with the map is denoted by \( \mathcal{F} \), which satisfies the universal property that for any sheaf \( \mathcal{G} \) and map of presheaves \( \psi : \presh \to \mathcal{G} \), there is a unique map \( \psi^\sharp : \mathcal{F} \to \mathcal{G} \) s.t. 
    \( \psi^\sharp \circ \text{sh} = \psi \).
    Without going full abstract nonsense (if we haven't done so already),
    it resolves the obstruction of gluing our local sections to a global section. This introduces objects which we will refer to as local risk measures.
    \begin{definition} \label{deflocalrm}
        A local risk measure \( \rho : U \to \extreals \) with open \( U \subseteq \mathcal{X} \) is a map 
        s.t. for each \( X \in U \), there is some open \( V \subseteq U \) on which \( \rho \)  restricts to an element in \( \extrm \).
    \end{definition}
    \begin{remark} \label{analoglocalrm}
        Definition \ref{deflocalrm} already looks similar to the sections of the sheaf from Proposition \ref{sheafification}
        in regard to local behavior. That they in fact specify the same object is however not immediate.
        First one needs to define the local analogues \( \extrmloc, \rejidbloc{U}, \accidbloc{U} \) verify 
        that \( \extrmloc \) is indeed a distributive bounded lattice which is in addition complete, JID/MID
        and that \( \rejidbloc{U} \) is a filter. This task is left to the reader, if they wish to do so.
        The actual isomorphism is something we will cover later.
    \end{remark}
    %(TODO it is not immediate how such objects can be used to match sections in the sheaf from \ref{sheafification}) \newline
    As concrete example why we need such objects for a sheaf and a counterexample that \( \presh \) from Theorem \ref{rmpresheaf}
    is only a presheaf, take any disjoint open sets \( U_1, U_2 \) in \( \mathcal{X} \) s.t. \( \rho_1 = \text{VaR}_{\lambda} > 0 \) on \( U_1 \) and 
    \( \rho_2 = \mathbb{E}[-X] > 0 \) on \( U_2 \), so \( \rho_1 \in \rejectidb{U_1}, \rho_2 \in \rejectidb{U_2} \), since \( U_1 \cap U_2 = \emptyset \) and \( \rho_1, \rho_2 \) 
    restrict to \( \rho_{\top} \) on \( \rejectidb{\emptyset} \), there must be a risk measure \( \rho \in \rejectidb{U_1 \cup U_2} \), if \( \presh \) would be a sheaf, 
    restricting to \( \rho_i \) on \( U_i \), but this seems impossible for non trivial \( \Omega, \mathbb{P} \). There are also more realistic examples:
    \begin{example}
        Let \( \rho_1 = \text{ess sup} - X \), the worst case risk measure. We will define some intermediate 
        risk measures for the actual counter example:
        \begin{align*}
            \eta_i(X) = \begin{cases}
                \text{ess sup} (- X) + i\cdot 100, \text{ for } i \leq 2 \\
                 \expv{-X} + i \cdot 1000, \text{ else}
            \end{cases}
        \end{align*}
        we combine them to risk measures 
        \begin{align*}
            \rho_2 &= \rho_1 \wedge (\eta_1 \vee \eta_5) \\
            \rho_3 &= (\rho_1 \vee \eta_3) \wedge (\eta_2 \vee \eta_4)
        \end{align*}
        with corresponding acceptance sets
        \begin{align*}
            \mathcal{A}_{\eta_i} = \begin{cases}
                \{ X \in \mathcal{X} | X \geq i \cdot 100 \text{ a.s. } \}, \text{ for } i \leq 2 \\
                \{ X \in \mathcal{X} | \expv{X} \geq i \cdot 1000 \}, \text{ else}
            \end{cases}
        \end{align*}
        and \( \mathcal{A}_{\rho_2}, \mathcal{A}_{\rho_3} \) the corresponding intersections and unions of them (recall \ref{acceptanceoperations}):
        \begin{align*}
            \mathcal{A}_{\rho_2} &= \mathcal{A}_{\rho_1} \cup (\mathcal{A}_{\eta_1} \cap \mathcal{A}_{\eta_5}) \\
            \mathcal{A}_{\rho_3} &= \left(\mathcal{A}_{\rho_1} \cap \mathcal{A}_{\eta_3} \right) \cup \left(\mathcal{A}_{\eta_2} \cap \mathcal{A}_{\eta_4} \right)
        \end{align*}

        Let \( U := B_{100}(1) \) then all \( X \in U \) are naturally in \( \mathcal{A}_{\rho_1} \) 
        and then also in \( \mathcal{A}_{\rho_2} \). By assumption \( 99 < \expv{X} < 101  \)
        so \( X \notin \mathcal{A}_{\eta_5} \), therefore \( \rho_2 \) coincides with \( \rho_1 \) on \( U \).
        Further \( X \notin \mathcal{A}_{\eta_3} \) and \( X \notin \mathcal{A}_{\eta_4} \), and therefore \( X \notin \mathcal{A}_{\rho_3} \)
        implying \( \rho_3 \in \rejectidb{U} \).
        Observe, that \( \rho_1 \wedge \eta_1 = \rho_1 \), and therefore \( \rho_1 \wedge (\eta_1 \vee \eta_5) = \rho_1 \).
        This enables the following set of equivalences:
        \begin{align*}
            \rho_1 \wedge (\eta_2 \vee \eta_4) &= \rho_1 \wedge (\eta_2 \vee \eta_4) \\
            \rho_1 \wedge (\rho_1 \vee \eta_3) \wedge (\eta_2 \vee \eta_4) &= \rho_1 \wedge (\eta_1 \vee \eta_5) \wedge (\eta_2 \vee \eta_4) \\
            \rho_1 \wedge \rho_3 &= \rho_1 \wedge (\rho_1 \vee \eta_3) \wedge (\eta_1 \vee \eta_5) \wedge (\eta_2 \vee \eta_4) \\
            \rho_1 \wedge \rho_3 &= \rho_2 \wedge \rho_3
        \end{align*}
        In particular, this means that \( [\rho_2] = [\rho_1] \) in \( \presh(U) \), but they can't be extended to a global risk measure
        in \( \extrm = \presh(\mathcal{X}) \), since they are distinct.
        The intuition behind the individual risk measures is that the viewpoint \( \rho_2 \) is also content with risky positions which have a high return 
        as opposed to \( \rho_1 \) which only takes on safe positions. \( \rho_3 \) requires a higher return then \( \rho_1 \) on safe positions, but is more lenient
        than \( \rho_2 \) for risky ones.
    \end{example}

    With the Definitions from Remark \ref{analoglocalrm} we now show the direct approach to our sheaf.
    \begin{theorem}
        There is an isomorphism \( \Psi : \extrmloc / \rejidbloc{U} \to \mathcal{F}(U) \) of bounded distributive lattices
        for each \( U \subseteq \mathcal{X} \).
    \end{theorem}
    \begin{proof}
        As an intuition of what is to come, one might refresh their repertoire on partitions of unity.
        %We will utilize a common trick in the construction of sheafs, which is conceptually the same as a partition of unity.
        Let \( \psi \in \extrmloc \) be a local risk measure and define \( \Psi([\psi]) \) as the section \( (\sigma_X)_{X \in U} \in \mathcal{F}(U) \)
        s.t. for each \( X \in U \) and corresponding neighborhood \( V \subseteq \) of \( X \) on which \( \psi \) restricts to a standard risk measure \( \rho \in \extrm \)
        we have \( \sigma_Y = \varphi_{\{Y\}, V}([\rho]) \). This already fullfills the \( (\ast) \) condition of Proposition \ref{sheafification}.
        \newline
        \newline
        We will show well-definedness of \( \Psi \) with injectivity at once.
        Take \( \psi_1, \psi_2 \in \extrmloc \) and assume \( \Psi([\psi_1]) = \Psi([\psi_2]) \) in \( \mathcal{F}(U) \),
        implying for each \( Z \in U \) and corresponding neighborhood \( V \) and restrictions \( \rho_1, \rho_2 \)
        of \( \psi_1, \psi_2 \) some \( \gamma_Z \in \rejectidb{V_Z} \) s.t. \( \gamma_Z \wedge \psi_1 = \gamma_Z \wedge \psi_2 \)
        and we can construct an open cover \( \{V_Z\}_{Z \in U} \) of \( U \) with them.
        As \( \mathcal{X} = L^\infty(\Omega, \mathcal{F}, \mathbb{P}) \) is a metric space, so is every open subset of it, 
        whence every \( U \) is paracompact yielding some locally finite refinement \( \{W_Z\}_{Z \in U'} \) of \( \{V_Z\}_{Z \in U} \).
        Define 
        \begin{align} \label{patchedlocalrm}
            \gamma(X) := \inf_{\{Z | X \in W_Z \}} \gamma_Z(X).
        \end{align}
        since \( X \) is covered by at least one \( W_Z \) the set over which we take infimum is non-empty 
        and for every included \( W_Z \) it holds \( \gamma_Z(X) > 0 \). Since \( \{W_Z\}_{Z \in U'} \) is locally finite,
        only finitely many \( W_Z \) are taken into account and therefore this is even a minimum, so \( \gamma(X) > 0 \).
        Since the set of \( W_Z \) over which we take the infimum is constant over some neighborhood of \( X \), by Definition
        of local finiteness, \( \gamma \) is a risk measure there and by arbitrariness of \( X \), it follows \( \gamma \in \rejidbloc{U} \).
        \iffalse
        Since for every \( X \in U \) there is some open neighborhood \( V \), which intersects only finitely many \( W_Z \), call them active, 
        which is again an open, we can focus on the intersection of \( V \) and all active \( W_Z \), which by construction is non-empty. 
        Here \( \gamma \) restricts just to a regular risk measure, namely \( \bigwedge_{Z^\ast} \gamma_{Z^\ast}(X) \), where \( Z^\ast \)
        correspond to the active \( W_Z \). As we could do this for each \( X \in U \), we have .
        \fi
        Moreover, we have 
        \[
            \psi_1 \wedge \gamma = \psi_1 \wedge \left (\bigwedge_{Z \in U'} \gamma_{Z} \right) = \bigwedge_{Z \in U'} \psi_1 \wedge \gamma_{Z}
        \]
        and in any specific \( W_Z \) it holds \( \psi_1 \wedge \gamma_Z = \psi_2 \wedge \gamma_Z \). Therefore the term rewrites to
        \[
            \bigwedge_{Z \in U'} \psi_2 \wedge \gamma_{Z} = \psi_2 \wedge \gamma,
        \]
        yielding \( [\psi_1] = [\psi_2] \) in \( \extrmloc / \rejidbloc{U} \).
        \newline
        \newline
        Let \( (\sigma_X)_{X \in U} \in \mathcal{F}(U) \) and extract a locally finite cover \( \{W_Z\}_{Z \in U'}\) from the cover induced by the defining \( V \)
        of the \( (\ast) \) property. 
        %On each \( W_Z \) we get at least one representative \( \rho \in \extrm \) with \( \sigma_Y = \varphi_{\{Y\}, V}([\rho]) \)
        %for \( Y \in W_Z \).
        Use the axiom of choice to pick a representative from each class \( [\rho_Z] \in \mathcal{F}(W_Z) \) with \( [\rho_Z] = \sigma |_{W_Z} \) and define
        \[
            \psi(X) := \inf_{\{Z | X \in W_Z\}} \rho_Z(X),
        \]
        deduce analogously to the construction of Equation \ref{patchedlocalrm}, that \( \psi \in \extrmloc / \rejidbloc{U} \).
        Suppose \( \sigma|_{V_X} = [\rho] \) with \( \rho \in \rejectidb{V_X} \), we then have for each \( Y \in V_X \) either \( \rho(Y) \leq 0 \)
        or \( \rho(Y) > 0 \). In the first case, also \( \rho_Z(X) = \rho(Y) \) for all \( \rho_Z \) with \( W_Z \cap V_X \neq \emptyset \), 
        since they are in the same equivalence class, therefore \( \psi(Y) = \rho(Y) \) and \( [\psi]_{Y} = \sigma_Y \). Otherwise all \( \rho_Z(X) > 0 \),
        but not necessarily with the same value. Regardless, it follows \( [\psi]_{Y} = [\top] = \sigma_Y \), and therefore \( [\psi]|_{V_X} = \sigma \mid_{V_X} \).
        By arbitrariness of \( X \), we conclude with \( \Psi(\psi) = \sigma \).
    \end{proof}

    \begin{corollary}
        There exists a sheaf \( \mathcal{F} : \textbf{Open}(\mathcal{X}) \to \textbf{BDLat} \) on \( \mathcal{X} \),
        which we call the rejection risk sheaf of \( \mathcal{X} \).
    \end{corollary}
    \begin{proof}
        This is just a restatement of Proposition \ref{sheafification}.
    \end{proof}
    %As noted before the presheaf \( \presh \) and \( \mathcal{F} \) produce for general \( U \subseteq \mathcal{X} \) not the same 
    Local risk measures are interesting in their own right, since they could reflect the fact that locally the behavior of an actor
    appear rational or consistent, but globally they contradict themselves. Meaning they are willing to take on unnecessary risk in some situations.
    But there's a special case, where we can guarantee a clear risk profile globally.
    %There's a special case
    %with the 
    %There is a special case when local risk measures are forced to be a global risk measure.
    \begin{proposition}
        Let \( U \subseteq \mathcal{X} \) be convex, then any local risk measure \( \psi : U \to \extreals \) coincides with a global risk measure
        on \( U \).
    \end{proposition}
    \begin{proof}
        Define for each \( X \in U \), the function \( f(m) = \psi(X + m) + m \) and let \( V \) be the corresponding neighborhood of \( X \)
        in which \( \phi \) restricts to some member of \( \extrm \). In \( V \) we have \( \partial_m f(m) = 0 \) therefore \( \psi \) is locally constant.
        As any convex \( U \) is connected, a locally constant function is globally constant, hence \( \psi(X + m) = \psi(X) - m \) for all \( m \in \mathbb{R} \)
        s.t. \( X + m \in U \).
        \newline
        \newline
        Let \( X, Y \in U \) s.t. \( X \leq Y \) and by definition the path \( Z_\lambda := \lambda Y + (1 - \lambda)X \) is contained in \( U \) for all \( \lambda \in [0, 1] \).
        For fixed \( \omega \in \Omega \) we have \( \partial_\lambda Z_\lambda(\omega) = Y(\omega) - X(\omega) \geq 0 \), so \( Z_\lambda \) is a.s. inreasing.
        Let \( \{ V_\lambda \}_{\lambda \in [0, 1]} \) be an open covering of the path \( Z_\lambda \) in its subspace topology with the \( V_\lambda \) 
        corresponding to the restriction neighborhoods of \( \psi \). Since the map \( \lambda \mapsto Z_\lambda \) is a.s. continuous, we obtain some open cover
        \( \{I_\lambda\}_{\lambda \in [0, 1]} \) of \( [0, 1] \) via its inverse. By compactness of \( [0, 1] \) this has some finite subcover \( \{I_\lambda\}_{\lambda \in \Lambda} \)
        and in each interval \( Z|_{I_\lambda} \) it holds for \( \lambda_1 \leq \lambda_2 \), \( \varphi(Z_{\lambda_1}) \geq \varphi(Z_{\lambda_2}) \). Since 
        the restrictions of \( \varphi \) coincide on overlaps of the \( Z_{I_\lambda} \), we obtain \( \varphi(Z_0) \geq \varphi(Z_1) \), so \( \varphi \) is a global risk measure
        on \( U \).
    \end{proof}

    In the following sections, we will discuss the meaning of robust representations from a lattice point of view and 
    by doing so stay very close to the underlying geometric interpretation. However, our sheaf construction won't re-enter the stage anymore. 
    The reader might wonder what they takeaway of this should be then. First, there are not many existing sheafs in literature 
    which map into the category of lattices, so this is a new addition to this collection. Secondly, it shows that modern tools of algebraic geometry
    can be translated to non-polynomial settings aside from the usual algebraic number theory. The obstruction to study the geometry of acceptance and rejection sets, 
    here is essentially that the ideals and filters are a very one-sided perspective. In algebraic geometry we have the luxury to be able to still evaluate sections on the underlying space,
    while for our quotient lattices this was only possible for the elements, which were accepted. Essentially one would need to consider two sheafs, 
    one for the rejection filters and another for acceptance ideals (which have not built here, but the motivated reader can translate the previous proofs to this setting) 
    to get the full picture, but the author perceived this as adding more unnecessary complexity than insight. A better behaved version might be the sheaf of local risk measures
    that is \( U \mapsto \extrmloc(U) \), where \( \extrmloc(U) = \{ \psi|_U | \psi \in \extrmloc \} \), but this is out of the scope of this thesis.
\end{subsection}

\end{section}
\newpage
\appendix
\begin{section}{Appendix A}
  %\begin{subsection}{Lattice Topologies}
    As stated before we loose Lipschitz continuity, if we complete the risk measure lattice,
    since expressions such as \( | \rho_{\top}(X) - \rho_{\top}(Y) | \) are undefined.
    Generally risk measures are, as far as the author is informed, not typically studied under the aspect
    of their own topology. 
    \newline
    The extended real numbers \( \extreals \) are second countable, compact, metrizable and \( \mathbb{R} \) is dense in \( \extreals \) 
    (Example 2.75 in \cite{infinitedimanalysis}).
    One way (there are more) to grant them a metric is by noting that there is a homeomorphism 
    \[\psi : \extreals \to [0, 1], x \mapsto \frac{x}{1 + |x|} \]
    and then to define \( d_{\extreals}(x, y) = d_{\mathbb{R}}(\psi(x), \psi(y)) = |\psi(x) - \psi(y)| \).
    Similarly we can proceed to endow \( \extrm \) with a metric.
    \begin{theorem} \label{unifconvrm}
        One can define a metric on \( \extrm \) via
        \[
        \dextrm{\rho_1}{\rho_2} = \sup_{X \in \mathcal{X}} d_{\extreals}(\rho_1(X), \rho_2(X)),
        \]
        where \( d_{\extreals}(\cdot, \cdot) \) is the metric of the extended real numbers.
        \( \extrm \) is with respect to the induced topology a complete space 
        in which \( \mathcal{R} \) is dense.
    \end{theorem}
    \begin{proof}
    Clearly \( \dextrm{\rho}{\rho} = 0 \) and \( \dextrm{\rho_1}{\rho_2} = \dextrm{\rho_2}{\rho_1} \), since \( d_{\extreals}(\cdot, \cdot) \) is a metric.
    Further if \( \dextrm{\rho_1}{\rho_2} = 0 \), then \( \rho_1(X) = \rho_2(X) \) for all \( X \in \mathcal{X} \) (as \( \psi \) is injective),
        and \( \rho_1 = \rho_2 \). Moreover,
        \begin{align*}
            \sup_{X \in \mathcal{X}} d_{\extreals}(\rho_1(X), \rho_3(X)) 
            &\leq \sup_{X \in \mathcal{X}} d_{\extreals}(\rho_1(X), \rho_2(X)) + d_{\extreals}(\rho_2(X), \rho_3(X)) \\
            &\leq \sup_{X \in \mathcal{X}} d_{\extreals}(\rho_1(X), \rho_2(X)) + \sup_{X \in \mathcal{X}} d_{\extreals}(\rho_2(X), \rho_3(X)) \\
            &= \dextrm{\rho_1}{\rho_2} + \dextrm{\rho_2}{\rho_3},
        \end{align*}
            where the triangle inequality of \( d_{\extreals}(\cdot, \cdot) \) is applied pointwise for the first inequality.
            Thus \( \dextrm{\cdot}{\cdot} \) is a metric.
            \newline
            Take a convergent sequence \( \{ \rho_n \}_{n \in \mathbb{N}} \subset \extrm \), 
            and any \( X, Y \in \mathcal{X} \) s.t. \( X \leq Y \), for each \( n \) it holds
            \begin{align*}
                \rho_n(Y) &\leq \rho_n(X) \\
                \psi(\rho_n(Y)) &\leq \psi(\rho_n(X)) \\
                \lim_{n \to \infty} \psi(\rho_n(Y)) &\leq \lim_{n \to \infty} \psi(\rho_n(X)) \\
                \psi(\lim_{n \to \infty} \rho_n(Y)) &\leq \psi(\lim_{n \to \infty} \rho_n(X)) \\
                \psi(\rho(Y)) &\leq \psi(\rho(X)) \\
                \rho(Y) &\leq \rho(X)
            \end{align*}
            due to \( \psi \) being monotone, continuous, invertible and non-strict inequalities being preserved by taking limits.
            As \( \rho(X + m) = \lim_{n \to \infty} \rho_n(X + m) = \lim_{n \to \infty} \rho_n(X) - m = \rho(X) - m \),
            \( \rho \in \extrm \), whence \( \extrm \) is complete.
            Consider the sequence \( \{ \rho_n \}_{n \in \mathbb{N}},\ \rho_n = \rho + n \) for some \( \rho \in \mathcal{R} \),
            clearly \( \lim_{n \to \infty} \rho_n = \rho_\top \) and similarly we can construct a sequence for \( \rho_\bot \)
            hence the closure of \( \mathcal{R} \) is \( \extrm \).
    \end{proof}
    As in the case of the continuous functions, the metric from Theorem \ref{unifconvrm} can be interpreted as being induced by uniform convergence
    of risk measures. As opposed to robustness in regard to the probability measure, risk measures which are close in this space, are in so far
    robust that it shouldn't matter much, which risk measure we use.
    The application of this analysis is that we can ensure that our lattice structure is compatible with a desirable topology 
    on \( \extrm \).
    \begin{lemma} \label{maxminarecont}
        The \( \max, \min : \extreals \times \extreals \to \extreals \) are continuous.
    \end{lemma}
    \begin{proof}
        We show the result for \( \max \), the proof for \( \min \) is similar.
        A subbasis for the topology of \( \extreals \) is given by all intervals of the form \( (x)^\upuparr \)
        and \( (x)^\dodoarr \) (see Example 2.2.6 in \cite{infinitedimanalysis}).
        We have \[ 
            (\max)^{-1} ((x)^\upuparr) = ((x)^\upuparr \times (x)^\downarrow) \cup ((x)^\downarrow \times (x)^\upuparr),
        \]
        and as each part of the union contains the open set \( (x)^\upuparr \) so is its product with any set
        and the union of two open sets is open as well.
    \end{proof}
    \begin{corollary}
        The lattice \( \extrm \) is a topological lattice.
    \end{corollary}
    \begin{proof}
        Since \( \extreals \) is compact, any continuous function on $\extreals \times \extreals$ is 
        uniformly continuous, in particular the \( \max, \min \) by Lemma \ref{maxminarecont}.
        Therefore, for every \(\epsilon > 0\), there exists a \(\delta > 0\) such that for any \(a, b, a', b' \in \extreals\):
        \[
        d_{\extreals}(a, b) < \delta \quad \text{and} \quad d_{\extreals}(a', b') < \delta 
        \]
        so by continuity \( d_{\extreals}(\max(a, a'), \max(b, b')) < \epsilon \).
        \newline

        Assume \(d_{\extrm}(\rho_{0}, \rho_{1}) < \delta\) and \(d_{\extrm}(\rho_0', \rho_1') < \delta\). 
        By the definition of \( \dextrm{\cdot}{\cdot} \), this implies that for all \(X \in \mathcal{X}\):
        \[
            d_{\extreals}(\rho_{0}(X), \rho_1(X)) < \delta \quad \text{and} \quad d_{\extreals}(\rho_0'(X), \rho_1'(X)) < \delta
        \]

        Using the uniform continuity from \( \extreals \), it follows that for all $X \in \mathcal{X}$:
        \[
            d_{\extreals}\big(\max(\rho_0(X), \rho_0'(X)), \max(\rho_1(X), \rho_1'(X))\big) < \epsilon.
        \]
        Taking the supremum over all \( X \in \mathcal{X}\):
        \[
            d_{\extrm}(\rho_0 \vee \rho_0', \rho_1 \vee \rho_1') \le \epsilon,
        \]
        and the result holds vis-à-vis for \( \wedge \).
    \end{proof}

   
    Now we can also give some topological classification to some algebraic results from before.
    For example Lemma \ref{necessaryhilbert} now also requires that any ideal \( P \)  
    in \( \extrm \) for which it holds \( \acceptidealb{\riskidealb{P}} = P \), is a closed convex set in \( \extrm \),
    meaning it is the intersection of the kernel of the linear functions \( f_X : \extrm \to \extreals, \rho \mapsto \rho(X) \)
    for \( X \in S \).
    %This could potentially be an analogue to the concept of convex ideals from real algebraic geometry (Definition \ref{defconvexideal}).
\end{subsection}

  %\begin{lemma}[Hahn-Banach alike for Riesz spaces] \label{hahnbanachriesz}
    For compact convex disjoint \( A, B \subset \mathcal{X} \), \( \mathcal{X} \) a Riesz-space, which can not be ordered relative to each other,
    meaning there are no \( a \in A \), \( b \in B \) with \( a \leq b \) or \( b \leq a \), there is a continuous linear
    order preserving \( f : \mathcal{X} \to \mathbb{R} \) and \( s, t \in \mathbb{R} \), s.t. 
    \begin{align*}
        f(a) < s < t < f(b)
    \end{align*}
    for all \( a \in A \), \( b \in B \). If in addition \( \mathcal{X} \) is a Banach algebra and has a unit then \( f(1) \neq 0 \).
\end{lemma}
\begin{proof}
    Let \( K = B - A \) be the Minkowski difference of \( A \) and \( B \), which is also convex and compact.
    Since \( A, B \) are not ordered relative to each other we have \( b - a \notin - \mathcal{X}_{+} \) or put differently
\( K \cap (- \mathcal{X}_{+}) = \emptyset \). As \( - \mathcal{X}_{+} \) is a convex linear cone, which is closed (\( X_{+} = \{X \in \mathcal{X} | X \geq 0 \} \)), we can apply the standard Hahn-Banach separation theorem to \( K \) and \( \mathcal{X}_{+} \) to find a continuous linear 
    \( f : \mathcal{X} \to \mathbb{R} \) and \( s, t \in \mathbb{R} \) s.t.
    \[
        f(x) < s < t < f(k)
    \]
    for all \( k \in K \) and \( x \in - \mathcal{X}_{+} \). Further \( f(x) \leq 0 \) (so it is order-preserving) 
    for all \( x \in - \mathcal{X}_{+} \), because if \( f(x) > 0 \) for some \( x \), 
    then we can find some \( \lambda \in \mathbb{R}_{+} \), s.t. \( f(\lambda x) > s \) and as \( - \mathcal{X}_{+} \) is a linear cone, 
    \( \lambda x \in - \mathcal{X}_{+} \).  Hence \( 0 < s \) and also 
    every \( k \in K \) can be written as \( k = b - a \), so we have 
    \begin{align*}
        f(k) &> 0 \\
        f(b - a) &> 0 \\
        f(b) &> f(a)
    \end{align*}
    by linearity. Since for any distinct real numbers there are infinitely many numbers in between them,
    we can find \( s, t \) as desired from \( \inf_{b \in B} f(b) > \sup_{a \in A} f(a) \).
    In case that \( \mathcal{X} \) is a Banach algebra, \( 0 < 1 \) and \( f(x) \leq 0 \) on \( \mathcal{X}_{+} \),
    where the equality is by linearity also achieved, imply that \( f(1) > 0 \) by order-preservation. 
\end{proof}

This Lemma enables us to derive a fact that comes often in handy to stipulate the existence of a risk measure with some restrictions on the acceptance set.
\begin{theorem} \label{riskmeasureseparatepoints}
    For every \( X, Y \in \mathcal{X} = L^p(\Omega, \mathcal{F}, \mathbb{P}) \) for \( 1 \leq p \leq \infty \) with 
    neither \( X \leq Y \) nor \( Y \leq X \) is satisfied, we can 
    find \( \rho_1, \rho_2 \in \mathcal{R} \) s.t. 
    \[ \rho_1(X) < 0, \quad \rho_1(Y) > 0 \]
    \[ \rho_2(X) > 0, \quad \rho_2(Y) < 0. \]
\end{theorem}
\begin{proof}
    \( \{X\}, \{Y\} \) are closed, trivially convex and even compact
    so by the order version of the Hahn-Banach separation from Theorem \ref{hahnbanachriesz},
    we can find some order preserving continuous linear \( \ell_1' : \mathcal{X} \to \mathbb{R} \) and \( \beta \in \mathbb{R} \), 
    s.t. 
    \[ \ell_1'(Y) < \alpha < \ell_1'(X) \]
    and by changing the roles of \( X \) and \( Y \) also a 
    \( \ell_2' : \mathcal{X} \to \mathbb{R} \) with
    \[ \ell_2'(X) < \beta < \ell_2'(Y) \]
    and both satisfy \( \ell_i(1) > 0 \).
    "Normalizing" \( \ell_i' \) with \( \ell_i := \frac{\ell_i'}{\ell_i'(1)} \) also yields a linear cont. monotone functional,
    which in addition satisfies for real constants \( c \in \mathbb{R} \), 
    \[ \ell_i(X + c) =  \ell_i(X) + \ell_i(c\cdot 1) = \ell_i(X) + c \cdot\ell_i(1) = \ell_i(X) + c. \]
    As \( \ell_i \) was order monotone, \( -\ell_i(X) = \ell_i(-X) \) is order antitone
    and therefore a coherent risk measure and \( \ell_i(-X) + c \) is still a monetary risk measure.
    Thus we can construct the risk measures 
    \begin{align*}
        \rho_1(Z) = \ell_1(-Z) + \frac{\alpha}{\ell_1(1)} \\
        \rho_2(Z) = \ell_2(-Z) + \frac{\beta}{\ell_2(1)}.
    \end{align*}
    From which we deduce,
    \begin{align*}
        \rho_1(Y) = \ell_1(- Y) + \frac{\alpha}{\ell_1(1)} > 0 > \ell_1(- X) + \frac{\alpha}{\ell_1(1)} \\
        \rho_2(X) = \ell_2(- X) + + \frac{\beta}{\ell_2(1)} > 0 > \ell_2(-Y) + \frac{\beta}{\ell_2(1)}.
    \end{align*}
\end{proof}
\begin{remark} \label{separationremark}
    Let \( X < Y \) then it is also no issue to find a risk measures that accepts \( Y \) but not \( X \).
    Take any risk measure \( \rho \in \mathcal{R} \), then define \( \rho'(Z) = \rho(Z) - c \) with \( c = \rho(Y) \in \mathbb{R} \)
    and \( \rho' \in \mathcal{R} \). It holds \( 0 = \rho'(Y) < \rho'(X) \). 
    However, \( \mathcal{R} \) has only antitone members, so there's no risk measure that accepts
    \( X \), but not \( Y \).
    Combined with Theorem \ref{riskmeasureseparatepoints} one can express this as \( \mathcal{R} \) separating points by acceptance,
    when \( \mathcal{X} \) is a Banach lattice. 
    The same results hold also for \( \extrm \) since the proofs are not dependent on any boundedness properties in regard to the risk measures.
\end{remark}

  \begin{subsection}{Colimits in Filter Quotient Lattices}
    The colimit has multiple definitions which are mostly equivalent, however
    the most useful one in this context is the categorical one with a universal property, because
    it automatically ensures that the objects in question are in the same category and we can make sense
    of operations in the colimit.
    \begin{definition} \label{indexcategory}
        Let \( \mathcal{I} \) be a small category (the objects and morphisms are sets), 
        which is filtered, meaning
        \begin{enumerate}
            \item for each \( x, y \in \mathcal{I} \) there is a \( z \in \mathcal{I} \) and morphisms \( x \to z \), \( y \to z \),
            \item for distinct morphisms \( \alpha : x \to y \), \( \beta : x \to y \) there is a morphism \( \gamma : y \to z \) 
                s.t. \( \gamma \circ \alpha = \gamma \circ \beta \),
        \end{enumerate}
        then we call \( \mathcal{I} \) an index category.
    \end{definition}
    Note that for a topological space \( \mathcal{X} \) the category of open sets \( U \subseteq \mathcal{X} \) and 
    morphisms being the inclusions, is an index category.
    Given a functor \( \mathcal{F} : \mathcal{I} \to \mathcal{C} \) into some category \( \mathcal{C} \), we
    can associate pictorially a diagram with the image of \( \mathcal{F}(\mathcal{I}) \) (alternatively just a subcategory),
    which is called a diagram indexed by \( \mathcal{I} \)
    \begin{definition} \label{defcolimit}
        The colimit of a diagram in  \( \mathcal{C} \) indexed by \( \mathcal{I} \) via a functor \( \mathcal{F} \),
        is an object \( \varinjlim_{\mathcal{I}} C_i \) along with morphisms \( \{\varphi_{k}\} \) denoted by 
        \( (\varinjlim_{\mathcal{I}} C_i, \varphi_k) \),
        s.t. for each object \( i \in \mathcal{I} \) there is a morphism \( \varphi_i : \mathcal{F}(i) \to \varinjlim_{\mathcal{I}} C_i \)
        and for each pair of objects \( i, j \in \mathcal{I} \) and morphism \( f_{i,j} \in \text{Mor}_{\mathcal{I}}(i, j) \), 
        it holds \( \varphi_j \circ \mathcal{F}(f_{i,j}) = \varphi_i \).
        Moreover \( \varinjlim_{\mathcal{I}} C_i  \) is universal in these requirements, meaning for any other \( T \in \mathcal{C} \) 
        satisfying the above, there is a unique \( \theta : \varinjlim_{\mathcal{I}} C_i \to T \) s.t. \( \theta \) commutes with the
        the morphisms of \( T \).
    \end{definition}
    We will also denote \( \mathcal{F}(f_{i,j}) \) with \( f_{i,j} \) for readability.
    The following Proposition is a verification which is a bit tedious. 
    The main idea is dependent on \( \mathcal{I} \) being an index category, which behaves similarly like the category of subsets of a set under inclusion, for which the analoguous result holds and how a filter contained in another projects almost like the identity in their quotient lattices.
    \begin{proposition} \label{colimitforquotientlattice}
        Let \( \mathcal{I} \) be an index category of filters in a bounded distributive lattice \( L \)
        with morphisms the inclusions as sets, then we can associate a functor \( \mathcal{F} \) to the category \( \quotlat{L} \) 
        of quotient lattices of \( \mathcal{L} \), sending each inclusion \( F_i \subseteq F_j \) to the morphism 
        \( f_{i,j} : \frac{L}{F_i} \to \frac{L}{F_j} \).
        The corresponding colimit satisfies 
        \[ 
            \varinjlim_{\mathcal{I}} \frac{L}{F_i} = \frac{L}{F},
        \]
        where \( F = \bigcup_{i \in \mathcal{I}} F_i \).
    \end{proposition}
    \begin{proof}
        By Remark \ref{dualfilterquotient} and more specifically Proposition \ref{quotientprojectionlattice}
        we have a well defined projection \( f_{i,j} : \frac{L}{F_i} \to \frac{L}{F_j}, [a]_{F_i} \mapsto [a]_{F_j} \) for \( F_i \subseteq F_j \),
        so the (covariant) functor exists.
        Also any such index category \( \mathcal{I} \) satisfies the second property \ref{indexcategory} vacously, 
        since between any two filters there can be at most one inclusion.
        \newline
        In general the union of filters in a lattice is not a filter, since the meet of elements in distinct filters
        are not necessarily in the union anymore, but \( \mathcal{I} \) being an index category circumvents this.
        Namely, let \( a \in F_i, b \in F_j  \) as \( \mathcal{I} \) is indexed, there is some 
        \( F_k \) s.t. \( F_j, F_i \subseteq F_k \) and as \( F_k \) is a filter it follows that \( a \wedge b \in F_j \).
        Hence the maps \( \varphi_i : \mathcal{F}(F_i) \to \frac{L}{F} \) are indeed lattice homomorphisms 
        and that they commute with the inclusions \( f_{i,j} : \mathcal{F}(F_i) \to \mathcal{F}(F_j) \) is by design
        (again Proposition \ref{quotientprojectionlattice}).
        We will also need the following observation:
        \begin{lemma} \label{utillemmacolim}
            For any \( a, b \in L \) s.t. \( [a]_F, [b]_F \in L / F \) are distinct, \( [a]_{F_i}, [b]_{F_i} \)
            are distinct for all \( F_i \).
        \end{lemma}
        \begin{proof}
             Assume the contrary, then there's some \( c \in F_i \) s.t. \( c \wedge b = c \wedge a \)
            but \( F_i \subseteq F \) so \( c \in F \) implying \( [a]_F = [b]_F \).
        \end{proof}
        The universal property is then verified as follows: Given another lattice quotient \( L / T \)
        with morphisms \( \psi_i : L / F_i \to L / T \) for all \( F_i \in \mathcal{I} \), define for \( a \in L \) and \( [a]_F \in \frac{L}{F} \),
        \( \theta([a]_F) := \psi_i([a]_{F_i}) \) for any \( F_i \). This is indeed independent of the choice of \( F_i \), let \( F_j \)
        be another object in \( \mathcal{I} \), since \( \mathcal{I} \) is filtered, there is some \( F_k \), s.t. \( F_i, F_j \subset F_k \)
        and as \( (L / T, \psi_i) \) must satisfy \( \psi_k \circ f_{i,k} = \psi_i \) and similar for \( F_i \), we have
        \begin{align*}
            \psi_i([a]_{F_i}) &= \psi_k \circ f_{i,k}([a]_{F_i}) = \psi_{k}([a]_{F_k}) \\
            \psi_j([a]_{F_j}) &= \psi_k \circ f_{j,k}([a]_{F_i}) = \psi_{k}([a]_{F_k}),
        \end{align*}
        hence \( \psi_i([a]_{F_i}) = \psi_j([a]_{F_j}) \). Note that this argument goes exactly the same to show that 
        \( \phi_j([a]_{F_i}) = \phi_i([a]_{F_j}) \) for any \( F_i, F_j \).
        \newline
        Take \( a, b \in L \) s.t. \( [a]_F = [b]_F \), then there is \( F_j \) with \( j \in F_j \), s.t. \( a \wedge j = b \wedge j \)
        (since any \( j \in F \) is in some \( F_j \)) and \( [a]_{F_j} = [b]_{F_j} \).
        By the preceeding it already holds \( \psi_i([a]_{F_i}) = \psi_j([b]_{F_j}) \) so 
        \[ 
            \theta([a]_{F}) = \psi_i([a]_{F_i}) = \psi_j([b]_{F_j}) = \theta([b]_{F}),
        \]
        hence \( \theta \) is a well defined map.
        \newline
        Take \( a, b \in L \) s.t. \( [a]_F \neq [b]_F \) then 
        \( \theta([a]_F \wedge [b]_F) = \theta([a \wedge b]_F) = \psi_i([a \wedge b]_{F_i}) \)
        and by Lemma \ref{utillemmacolim} \( [a \wedge b]_{F_i} = [a]_{F_i} \wedge [b]_{F_i} \), hence
        \( \psi_i([a \wedge b]_{F_i}) = \psi_i([a]_{F_i}) \wedge \psi_i([b]_{F_i}) = \theta([a]_{F}) \wedge \theta([b]_{F}) \).
        The argument for \( [a \vee b]_F \) is similar, so \( \theta \) is indeed a lattice homomorphism.
        \newline
        By construction it holds \( \theta \circ \phi_i([a]_{F_i}) = \theta([a]_{F}) = \psi_i(a_{F_i}) \) and for \( F_j \neq F_i \), 
        combine \( \phi_j([a]_{F_i}) = \phi_i([a]_{F_j}) \) and \( \psi_i(a_{F_i}) = \psi_j(a_{F_i}) \) to 
        \begin{align*}
            \theta \circ \phi_j([a]_{F_i}) = \theta \circ \phi_i([a]_{F_i}) = \psi_i([a]_{F_i}) = \psi_j(a_{F_i})
        \end{align*}
        Consequentially, \( \theta \circ \phi_j = \psi_i \) is true for all \( [a]_{F_j} \).
        \newline
        Suppose there is a homomorphism \( \tilde{\theta} : L / F \to L / T \) doing as desired and distinct from \( \theta \).
        Then there's some \( [a]_F \in L / F \), s.t. \( \tilde{\theta}([a]_F) \neq \theta([a]_F) \), 
        but \( [a]_F \) has a preimage \( [a]_{F_i} \) under \( \phi_i \),
        and then \( \tilde{\theta} \circ \phi_i ([a]_{F_i}) \neq \theta \circ \phi_i([a]_{F_i}) = \psi_{i}([a]_{F_i} \), 
        contradicting the hypothesis. 
    \end{proof}
    One might be able to argue in abstract nonsense, that the quotient lattices are subcategory of \( \textbf{DLat} \),
    which itself is also a variety in the sense of universal algebra and any such category is cocomplete, hence colimits exist \cite{univalgarecocomplete}.
\end{subsection}

\end{section}

\bibliographystyle{plain} % We choose the "plain" reference style
\bibliography{refs} % Entries are in the refs.bib file
\end{document}



